\documentclass[11pt]{article}
\usepackage[T1]{fontenc}
\usepackage{lmodern}
\usepackage[margin=1in]{geometry}
\usepackage{amsmath,amssymb,amsthm,mathtools}
\usepackage{booktabs,longtable,array,enumitem,microtype}
\usepackage{xcolor,listings,fancyhdr}
\usepackage[colorlinks=true,linkcolor=blue!45!black,urlcolor=blue!45!black,citecolor=blue!45!black]{hyperref}
\newtheorem{theorem}{Theorem}[section]
\newtheorem{lemma}[theorem]{Lemma}
\newtheorem{proposition}[theorem]{Proposition}
\newtheorem{corollary}[theorem]{Corollary}
\theoremstyle{definition}
\newtheorem{definition}[theorem]{Definition}
\newtheorem{remark}[theorem]{Remark}
\newcommand{\ipra}{\mathrm{iPRA}}
\newcommand{\Raw}{\mathcal R}
\newcommand{\Can}{\mathcal Z}
\newcommand{\Old}{\mathcal E}
\newcommand{\nf}{\mathop{\mathrm{nf}}}
\newcommand{\FV}{\mathop{\mathrm{FV}}}
\newcommand{\Eq}{\mathsf{Eq}}
\newcommand{\D}{\mathsf D}
\newcommand{\Q}{\mathsf Q}
\newcommand{\Use}{\mathsf U}
\newcommand{\Refl}{\mathsf R}
\newcommand{\Eval}{\mathsf V}
\newcommand{\cl}[2]{\langle #1;#2\rangle}
\newcommand{\llbracket}{[\![}
\newcommand{\rrbracket}{]\!]}
\DeclareRobustCommand{\file}[1]{\nolinkurl{#1}}
\setlist{itemsep=3pt,topsep=5pt}
\pagestyle{fancy}
\fancyhf{}
\fancyhead[L]{\small A smaller clausal presentation of iPRA}
\fancyhead[R]{\small\thepage}
\setlength{\headheight}{14pt}
\lstset{basicstyle=\ttfamily\scriptsize,breaklines=true,columns=fullflexible,
  keepspaces=true,showstringspaces=false,frame=single,rulecolor=\color{black!20},
  xleftmargin=3pt,xrightmargin=3pt}
\hypersetup{pdftitle={A smaller clausal presentation of iPRA},pdfauthor={Research report prepared for Vladimir Reshetnikov}}
\title{A smaller clausal presentation of iPRA\\[4pt]
  \large Root-free evaluation, a uniform companion, and checked proof translations}
\author{Prepared for Vladimir Reshetnikov}
\date{November 11, 2026}
\begin{document}
\maketitle
\begin{abstract}
We combine the repository's clause presentation with its elimination of the
direct equality-composition case. The resulting raw calculus has terms
$x,0,S(t),I(n,a;u.h)$, formulas $s=t$ and $\D_{x,y}B$, and five instruction
patterns: reflexivity, two-case evaluation, clause use, paired elimination,
and paired introduction. It has ordinary capture-avoiding substitution.
The complete standalone checker has 2,956 normalized bytes, 837 significant
tokens, and 977 AST nodes. Its matched control has 3,003 bytes, 842 tokens,
and 984 AST nodes; the exact incoming raw-clausal checker has 3,046 bytes.
The arithmetic deletion alone accounts for 47 bytes, separately from an
ordinary helper cleanup. Every open sequent is preserved literally. A
proof with $N$ nodes and $k$ direct composition nodes translates to a proof
with exactly $N+2k$ nodes, without duplicating its premise proofs.

A zero-normal companion uses one successor-evaluation schema and the same
five instruction patterns. A fixed-zero-base introduction guard removes the literal-witness test.
Its checker has 2,948 bytes, 852 tokens, and 987 AST nodes; it uses more
tokens and AST nodes than the raw proposal. Both ordinary realizations are
bi-interpretable and equiconsistent with full first-order intuitionistic
primitive recursive arithmetic. Both numerical domains are whole, equality
is actual equality, and the numerical composites have graph $x=y$.
The component mathematical reductions are inherited; the contribution is
their explicit integration, comparison, and executable proof translations.
No global simplicity minimum, new compact-codec record, or new Lean/Rocq
verification is claimed.
\end{abstract}
\tableofcontents

\section{Result and comparison target}
\label{sec:result}

The raw calculus $\Raw$ in Section~\ref{sec:raw-spec} is the recommended
profile when ordinary substitution and the token/AST coordinates matter most. It gives a strict improvement in all three
reported source-size coordinates over a matching implementation of the
current raw clausal calculus. The zero-normal calculus $\Can$ in Section~\ref{sec:canonical} has the
smallest byte count among the selected profiles, one arithmetic evaluation
schema, and a fixed-zero-base introduction test. Its normalizing substitution
is explicitly counted.
These are two complete programs, not three coordinatewise minima assembled
from unrelated programs.

The comparison is pinned to ProveIt commit
\begin{center}\small\texttt{831081033253cca2115196cf86bb99c77ed2aeae}.\end{center}
The reviewed subtree is \file{Logic/PrimitiveRecursiveArithmetic/ConciseIPRA}.
Its README describes fifty filed research packages and a Lean formalization.
Five additional ZIP archives were present in its \file{docs/incoming}
directory. In particular, the following results already existed:
\begin{itemize}
\item the Eq/$D$ clause interface, removing the separate $Q$ formula constructor;
\item three-case computational equality, which combines computation and transport;
\item removal of the direct equality-composition branch by a zero-iterator frame;
\item zero-normal terms with reflexivity and one successor-evaluation schema;
\item a paired-introduction guard based on the reflexivity of its zero instance.
\end{itemize}
Thus the present package does not claim priority for any of these component
ideas. Its raw calculus combines the first and third results; its canonical
companion combines the first and fourth. The source ledger in
\file{evidence/provenance.json} identifies the actual comparison material,
including the incoming archives. The main arithmetic bridge is included in
the appendices, with its upstream attribution.

The word ``simpler'' here has a precise, limited meaning: a smaller complete
standalone checker in the stated matching family, with a strictly reduced
set of arithmetic cases. There is no proved lower bound over all possible
presentations, languages, encodings, or host languages. The much smaller
compressed checker profiles already listed in the repository remain smaller.

\section{The complete raw axiomatization}
\label{sec:raw-spec}

\subsection{Numerical terms, formulas, and clauses}

There is one numerical sort. Terms are finite syntax trees
\[
 t ::= x\mid 0\mid S(t)\mid I(n,a;u.h).
\]
Only $u$ in $h$ is bound in an iterator. The counter $n$ and initializer $a$
are outside that binder. Bodies may contain fixed numerical parameters;
there are no variables over functions and no universal evaluator.
Substitution is ordinary capture-avoiding substitution, including under
iterator binders. Put $1=S0$.

Formulas and retained context clauses are
\[
 A ::= (s=t)\mid\D_{x,y}A,
 \qquad \cl{A}{y.B(y)},\quad y\notin\FV(A).
\]
Both $x$ and $y$ are bound in $\D_{x,y}A$, and are distinct. A context is
a finite list of clause occurrences. The clauses are proof-context data,
not additional objects in the numerical domain. Their ordinary readings are
\[
 (\D_{x,y}B)^\circ=\exists x\bigl(B^\circ(x,0)\to
                                  \forall y\,B^\circ(x,y)\bigr),
 \qquad
 \cl{A}{y.B}^\circ=A^\circ\to\forall y\,B^\circ.
\]
An ordinary formula assumption $F$ abbreviates the constant clause
$\cl{H}{v.F}$, where $H$ is the formula $0=0$ and $v$ is unused.
The designated contradiction is $0=S0$.

\subsection{All primitive rules}

The following display is the entire primitive inventory.
\paragraph{Reflexivity.}
\[
                 \frac{}{\Gamma\vdash t=t}\quad(\Refl).
\]
\paragraph{Evaluation.}
There are two arithmetic cases of one evaluator instruction:
\[
 \frac{\Gamma\vdash a=t\qquad
       \Gamma\vdash I(0,a;u.h)=r}
      {\Gamma\vdash r=t}\quad(\Eval_0),
\tag{V0}\label{rule:v0}
\]
\[
 \frac{\Gamma\vdash I(n,a;u.h)=t\qquad
       \Gamma\vdash I(Sn,a;u.h)=r}
      {\Gamma\vdash r=h[t/u]}\quad(\Eval_S).
\tag{VS}\label{rule:vs}
\]
The repeated terms in these patterns must match syntactically, up to a
fixed implementation of bound-variable renaming. In particular no test of
semantic equality is hidden in the side conditions. The body $h$ is read
from the second equality premise, not supplied as a separate transport
template.

\paragraph{Clause use.}
\[
 \frac{\Gamma\vdash A}{\Gamma\vdash B(t)}\quad(\Use),
 \qquad\cl{A}{y.B(y)}\in\Gamma.
\tag{U}\label{rule:u}
\]
\paragraph{Paired elimination.}
\[
 \frac{\Gamma\vdash\D_{x,y}B(x,y)\qquad
       \Gamma,\cl{B(b,0)}{y.B(b,y)}\vdash F}
      {\Gamma\vdash F}\quad(\mathsf E).
\tag{E}\label{rule:e}
\]
The eigenparameter $b$ is fresh for $\Gamma$, $F$, and all passive matrix
parameters.

\paragraph{Paired introduction.}
\[
 \frac{\Gamma,B(t,\eta)\vdash B(t,Sz)}
      {\Gamma\vdash\D_{x,y}B(x,y)}\quad(\mathsf I),
 \qquad
 \eta=\begin{cases}
 z,&B\text{ is a literal equation and }t\text{ is literally }0,\\
 0,&\text{otherwise}.
 \end{cases}
\tag{I}\label{rule:i}
\]
The eigenparameter $z$ is fresh for $\Gamma$, $t$, and all passive matrix
parameters. The displayed local formula assumption is its constant clause,
as defined above, and is discharged. An equation under a $D$ binder is not
a literal equation.

There are five instruction tags: $\Refl,\Eval,\Use,\mathsf E,\mathsf I$.
If each arithmetic row of $\Eval$ is counted separately, the display has
six schemas. The introduction also has an explicit mode test. These
conventions are stated to avoid obtaining an apparent reduction merely by
renaming several rules as one. The real deletion is the third arithmetic
case described next. No additional computation, equality, logical,
successor, induction, or structural axiom is primitive.

\section{Deleting the direct equality case}
\label{sec:deletion}

Let $\Old$ be the incoming raw clausal calculus. It has exactly the same
syntax and rules as $\Raw$, with one extra evaluator case
\[
       \frac{s=t\qquad s=r}{r=t}.\tag{VId}\label{rule:vid}
\]
All contexts in this section may contain arbitrary retained clauses.

\begin{lemma}[An iterator frame recovers direct composition]
\label{lem:frame}
Rule~\eqref{rule:vid} is derivable using only reflexivity and
\eqref{rule:v0}. Its two input proofs each occur once.
\end{lemma}
\begin{proof}
Given $p:s=t$ and $q:s=r$, put $W(s)=I(0,s;u.0)$. Reflexivity gives
$W(s)=W(s)$. Apply \eqref{rule:v0} first to $q$ and that reflexivity:
\[
                     W(s)=r.
\]
Apply it again to $p$ and this equation, giving $r=t$.
In proof notation the construction is
\[
         \Eval_0\bigl(p,\Eval_0(q,\Refl(W(s)))\bigr).
\]
No induction, clause manipulation, successor computation, numerical cases,
or prior symmetry principle is used.
\end{proof}

\begin{theorem}[Exact preservation of every open sequent]
\label{thm:raw-equivalence}
For every raw clause context $\Gamma$ and raw formula $A$,
\[
                 \Gamma\vdash_{\Old} A
                 \quad\Longleftrightarrow\quad
                 \Gamma\vdash_{\Raw} A.
\]
Both transformations are effective on finite proof trees. Numerical terms,
formulas, contexts, and free numerical parameters are unchanged.
\end{theorem}
\begin{proof}
Every $\Raw$ rule is an $\Old$ rule. In the other direction, induct on
the $\Old$ derivation and replace each use of \eqref{rule:vid} by
Lemma~\ref{lem:frame}. All other nodes are copied with their already
translated premise proofs. No variable is bound or freed by the replacement,
so all introduction and elimination freshness conditions remain valid.
\end{proof}

\begin{corollary}[Exact proof-node overhead]
\label{cor:overhead}
If the input tree has $N$ inference nodes and $k$ uses of
\eqref{rule:vid}, the displayed translation has exactly $N+2k$ nodes.
In particular its size in inference nodes is at most $3N$.
\end{corollary}
\begin{proof}
Each replaced node becomes two evaluation nodes and one new reflexivity
leaf. Both original premise subtrees are inserted once. All other nodes
are unchanged.
\end{proof}

This is not a bound on JSON byte length: the new frame contains a copy of
the source term. The shipped compiler preserves the end formula literally
and checks its result using the complete target checker. Its reverse
inclusion is the identity on proof syntax.

\section{Recovering arithmetic and logic}
\label{sec:recovery}

Theorem~\ref{thm:raw-equivalence} already transfers the earlier derivability
and interpretation results. We give the principal derivations explicitly
so the new inventory is independently understandable.

\subsection{Equality and the iterator equations}

Write $\operatorname{Root}(p,q)$ for Lemma~\ref{lem:frame}. Given
$p:s=t$, its symmetry is
$\operatorname{Root}(\Refl(s),p):t=s$. Given additionally $q:t=v$,
transitivity is $\operatorname{Root}(q,\operatorname{Sym}(p)):s=v$.
The two computation equations are evaluations of reflexivities:
\[
 \Eval_0(\Refl(a),\Refl(I(0,a;h))) : I(0,a;h)=a,
\]
\[
 \Eval_S(\Refl(I(n,a;h)),\Refl(I(Sn,a;h)))
 : I(Sn,a;h)=h[I(n,a;h)/u].
\]

\begin{lemma}[Congruence for any fixed term template]
\label{lem:congruence}
From $s=t$, derive $h(s)=h(t)$ for every finite template $h(u)$,
including multiple occurrences and occurrences under nested iterator binders.
\end{lemma}
\begin{proof}
Let $A_s=I(0,s;u.h)$ and $B_s=I(1,s;u.h)$. The following four evaluations
are all primitive cases, with no appeal to a derived transport rule:
\[
\begin{aligned}
 z_s&=\Eval_0(\Refl(s),\Refl(A_s))&&:A_s=s,\\
 c_s&=\Eval_S(z_s,\Refl(B_s))&&:B_s=h(s),\\
 z_{s,t}&=\Eval_0(p,\Refl(A_s))&&:A_s=t,\\
 \Eval_S(z_{s,t},c_s)&&&:h(s)=h(t).
\end{aligned}
\]
Substitution into the iterator body is precisely the primitive substitution
in \eqref{rule:vs}, so the argument applies to the full term grammar.
\end{proof}

Congruence followed by $\operatorname{Root}$ gives one-term crossed
transport, $s=t, h(s)=r\Rightarrow r=h(t)$. Two such transports give
transport in an arbitrary equation template.

\subsection{All numerical cases and equation induction}

Ordinary assumption use is derivable: prove $H$ by reflexivity and apply
U to its constant clause. Weakening, fresh renaming, and ordinary-assumption
cut are admissible by induction on derivations. For cut, each U that selects
the distinguished constant-clause occurrence is replaced by the supplied
proof; all other nodes are translated recursively after freshening
eigenparameters. These arguments do not assume that arbitrary substitution
preserves the introduction mode.

All case introductions are derivable:
\[
  \frac{\Gamma,B(t,0)\vdash B(t,Sz)}
       {\Gamma\vdash\D_{x,y}B(x,y)}.\tag{Case}
\]
Only an equation at the literal witness $0$ is missing from the primitive
inventory. For that instance use $c=I(0,0;u.0)$ as the witness. It is
syntactically different from $0$ and is provably equal to $0$.
Equation transport and ordinary cut move the assumed and resulting
instances between $c$ and $0$. Primitive introduction at $c$ is in case
mode and gives the required conclusion.

Arbitrary capture-avoiding raw substitution is now admissible. Induct on
the derivation after freshening eigenparameters. A progressive instance
keeps its equation root and literal-zero witness. A Case instance can
change mode only when its substituted witness becomes literal zero;
use the unrestricted derived Case rule in that event. All other rules
commute with substitution. Thus the ordinary logical macro proofs below
may substitute parameters without an unproved mode-stability assumption.

Transport through a compound formula follows by external induction on its
syntax. For $\D_{x,y}B(v,x,y)$, eliminate the old paired formula at a fresh
$b$, introduce the new one by Case at $b$, transport its assumed zero
instance backwards, use the opened old clause at $Sz$, and transport the
successor instance forwards. The recursive transports concern the strictly
smaller matrix. This establishes full formula transport without arithmetic
induction.

For an unused $w$, write $q_yB(y)=\D_{w,y}B(y)$. Case gives its introduction
from $B(0)\vdash B(Sz)$; E and U give its specialization from $q_yB$ and
$B(0)$ to $B(t)$. A retained clause $\cl{B(0)}{y.B(y)}$ itself proves
$q_yB$ by Case. These are the same three interfaces as the old native $Q$.

For a literal equation $E(y)$, primitive I with unused witness $0$ gives
$q_yE$ from $E(z)\vdash E(Sz)$. Therefore
\[
 \frac{\Gamma\vdash E(0)\qquad\Gamma,E(z)\vdash E(Sz)}
      {\Gamma\vdash E(t)}\tag{EI}
\]
is derived. Likewise Case, E, and U give numerical cases for arbitrary
formulas from their zero and successor branches. Progress has not been
granted to arbitrary quantified formulas.

\subsection{An explicit ordinary logical vocabulary}

The entire logical interface is recalled here; the detailed inherited
proofs are in Appendix~\ref{app:choice}. Define
$C(v,a,b)=I(v,a;u.b)$ with the bound $u$ unused. Then
$C(0,a,b)=a$ and $C(Sv,a,b)=b$.
Define a gate on formulas by
\[
 G_v(s=t)=(C(v,0,s)=C(v,0,t)),\qquad
 G_v(\D_{x,y}B)=\D_{x,y}G_vB,
\]
renaming bound variables away from the control. The formulas $G_0A$ are
theorems, and $G_{Sv}A$ is interderivable with $A$. Padding is
$A\dashv\vdash\D_{u,y}G_yA$, with $u,y$ fresh.

Select between equations by
\[
 M_v(s=t,r=w)=(C(v,s,r)=C(v,t,w)).
\]
Align paired binders recursively; if one argument is an equation and the
other is $D$-headed, pad the equation using $G_y$ under that binder.
Explicitly, with freshly aligned $x,y$,
\[
\begin{aligned}
 M_v(\D_{x,y}A,\D_{x,y}B)&=\D_{x,y}M_v(A,B),\\
 M_v(E,\D_{x,y}B)&=\D_{x,y}M_v(G_yE,B),\\
 M_v(\D_{x,y}A,E)&=\D_{x,y}M_v(A,G_yE).
\end{aligned}
\]
This recursion decreases the maximum binder depth. It gives
$M_0(A,B)\dashv\vdash A$ and $M_{Sv}(A,B)\dashv\vdash B$.
Put $\delta(v)=C(v,1,0)$ and $p(z)=I(z,0;u.\delta(u))$.
The following are finite macros, not extra primitive symbols:
\[
\begin{aligned}
 A\Rightarrow_*B&=\D_{u,z}M_z(A,B),\\
 A\wedge_*B&=\D_{u,z}G_zM_{p(z)}(A,B),\\
 \exists_*x\,A(x)&=\D_{x,y}G_yA(x),\\
 A\vee_*B&=\exists_*x\,M_x(A,B),\\
 \forall_*x\,A(x)&=A(0)\wedge_*\D_{u,z}G_zA(z),\\
 \bot_*&=(0=S0).
\end{aligned}
\tag{Logic}\label{eq:logic}
\]
Unused coordinates and all fresh binders are chosen capture-free.
Implication follows from the selector endpoints; the existential from
the gate; conjunction uses $p(1)=1$ and $p(2)=0$; disjunction uses
controls $0$ and $1$ and numerical cases. Universal elimination combines
the zero conjunct and the successor instances. Under $0=S0$, selector
transport gives any equation, and induction on formula syntax gives
explosion for every formula. These constructions establish intuitionistic
logic, without excluded middle for arbitrary formulas.

\section{A companion with one arithmetic evaluation schema}
\label{sec:canonical}

The raw system gives the cleanest exact comparison. The companion $\Can$
trades a normalizing substitution operation for one arithmetic schema and
a fixed-zero-base introduction test.

\subsection{Canonical syntax and substitution}

Let $\nf$ recursively contract every iterator whose counter is the literal
$0$:
\[
 \nf(I(n,a;u.h))=
 \begin{cases}
   \nf(a),&\nf(n)=0,\\
   I(\nf(n),\nf(a);u.\nf(h)),&\text{otherwise}.
 \end{cases}
\]
It acts structurally on variables, $0$, successor, formulas, and clauses.
Numerical terms in $\Can$ are exactly the fixed points of $\nf$.
Every substitution is followed by $\nf$; this includes contractions first
exposed by opening a binder. Successor-counter iterators are not reduced
by this normalizer.

Use the same clause rules and reflexivity as before, all with this canonical
substitution. Replace both raw evaluation rows by the single schema
\[
 \frac{\Gamma\vdash\nf(I(n,a;u.h))=t\qquad
       \Gamma\vdash I(Sn,a;u.h)=r}
      {\Gamma\vdash r=\nf(h[t/u])}.
\tag{UE}\label{rule:ue}
\]
There is no equality-composition case and no zero-evaluation case.
Replace the earlier introduction by
\[
 \frac{\Gamma,B(t,\eta)\vdash B(t,Sz)}
      {\Gamma\vdash\D_{x,y}B(x,y)},\qquad
 \eta=\begin{cases}
 z,&\nf(B(t,0))\text{ is literally }(0=0),\\
 0,&\text{otherwise}.
 \end{cases}
\tag{I0}\label{rule:i-zero}
\]
Freshness is unchanged. The test is one exact comparison of the computed
base with a fixed formula. It automatically excludes compound matrices,
because numerical substitution and zero normalization preserve formula
roots. There is no separate equation-root or witness test. The total
inventory is five instruction tags: Refl, UE, U, E, and I0. Zero computation
has moved into syntax; its implementation is counted in the checker.

\begin{lemma}[Canonical equality bootstrap]
\label{lem:canonical-equality}
Reflexivity and UE derive direct composition, every numerical congruence,
and successor computation in $\Can$.
\end{lemma}
\begin{proof}
For direct composition from $p:s=t$ and $q:s=r$, put
$K(s)=I(1,s;u.u)$. Since $\nf(I(0,s;u.u))=s$,
\[
    \operatorname{UE}\bigl(p,\operatorname{UE}(q,\Refl(K(s)))\bigr)
\]
proves $r=t$. The two UE instances have the same identity body.
Symmetry and transitivity follow as in the raw calculus.

For congruence put $J=I(1,s;u.h)$. UE on reflexivities gives
$J=\nf(h[s/u])$; UE on $p$ and that equation gives
$\nf(h[s/u])=\nf(h[t/u])$. Successor computation is UE on the two
appropriate reflexivities. Thus no transport rule is presupposed by the
evaluator's equality bootstrap.
\end{proof}

\subsection{A fixed zero-base test still recovers all induction}

Ordinary assumption use, weakening, and ordinary-assumption cut are proved
as before. Every Case instance is especially simple to recover. If its
base is not literally $0=0$, I0 is already Case. If the base is $0=0$,
cut its reflexivity proof into the supplied Case branch, then weaken that
branch by the progressive assumption and apply I0. No artificial witness
equal to zero is needed. Consequently all three $q$ interfaces and numerical
case reasoning are available before any equation induction is recovered.

Here is the full induction recovery, including the point where an arbitrary
equation is changed into one with the required base. Suppose
$E(y)$ is $s(y)=t(y)$ and a step proof
$\Gamma,E(z)\vdash E(Sz)$ is supplied. To introduce $q_yE(y)$ by Case,
retain $E(0)$ and aim to prove $E(Sw)$ for fresh $w$.
Put $s_0=s(0)$, $t_0=t(0)$ and
\[
 H(y):\quad C(y,0,s(y))=C(y,0,t(y)).
\]
In canonical syntax $H(0)$ is literally $0=0$. Consider the two decoders
\[
\begin{aligned}
 F_s(y):\quad C(y,s_0,C(y,0,s(y)))&=s(y),\\
 F_t(y):\quad C(y,s_0,C(y,0,t(y)))&=t(y).
\end{aligned}
\]
Their zero instances are respectively reflexivity and the retained $E(0)$.
Every successor instance is an equality obtained by two selector
computations. Thus Case gives each conditional closure, and its supplied
base lets the derived $q$-specialization (E followed by U) give
$F_s(z)$ and $F_t(z)$.

Under $H(z)$, congruence for $v\mapsto C(z,s_0,v)$, together with the two
decoder equations, gives $E(z)$. Cut this into the original step proof to
obtain $E(Sz)$, then use selector computation to infer $H(Sz)$. I0 is
therefore applicable to the $H$ matrix with an unused existential
coordinate: its zero instance is exactly $0=0$. It gives $q_yH(y)$.
Reflexivity supplies its base, and derived $q$-specialization at $Sw$
gives $H(Sw)$, hence $E(Sw)$
by computation. The outer Case now proves $q_yE(y)$.

This derives unrestricted equation progress. With a supplied base it derives
EI. The proof uses no equation induction before I0 is applied: the decoders
use Case, and their equality reasoning comes from
Lemma~\ref{lem:canonical-equality}. Conversely, when the guard is true, equation progress first gives
$q_yB(t,y)$. Its reflexive zero base and derived specialization give each
successor instance; Case at the original witness $t$ then gives the I0
conclusion. When the guard is false, I0 already is Case. Hence the fixed-base and separate Case/progress presentations
have exactly the same consequences. This guard and its symmetric lift are
inherited from the incoming reflexive-base report; the integration here
uses them with clausal single evaluation.

\subsection{Why normalization preserves the whole theory}

The necessary inherited normalization theorem says that the full native
calculus proves a raw sequent exactly when its zero-normal counterpart
proves the normalized sequent. The Lean source contains the corresponding
theorem \file{nzprf_iff_nprf}; the incoming single-evaluator report gives
the Refl/UE equality presentation used here.

There are two substantive points behind that theorem. First,
$\nf(\nf(t)\sigma)=\nf(t\sigma)$, and canonical substitution is compatible
with composition. Second, the ordinary equality $t=\nf(t)$ must also hold
when normalization changes an iterator's body. Given pointwise equal
bodies and equal seeds, equation induction on the iterator counter proves
equality of the two iterators: the base is seed equality, and the step uses
the supplied body equality and the induction hypothesis. External
induction on term syntax now proves $t=\nf(t)$, using ordinary term
congruence also to change the counter from $n$ to $\nf(n)$.

The exact-zero guard behaves well under substitution in one direction:
if the base is literally $0=0$, it stays so. A formerly nonzero base can
become $0=0$, changing a primitive Case instance into a progressive one.
That instance is repaired by the derived unrestricted Case rule above.
Thus arbitrary substitution is admissible, although keeping every original
rule tag is not a correct proof transformation. Equivalently, pass through
the calculus with separate Case and equation-progress rules and recover I0
afterwards. This is part of the mathematical specification, not a missing
checker side condition.

Finally replace native $Q_yB$ by $\D_{w,y}B$ with $w$ unused. In the other
direction, a clause introduced by E is represented by its $Q$ assumption,
whereas an ordinary constant clause is represented by its formula
assumption. U translates respectively to QE or Hyp. The three closure
interfaces from Section~\ref{sec:recovery} prove the forward direction.
Occurrence annotations distinguish clauses of different origins even if
their printed values coincide. These translations preserve normal terms
and the exact-zero guard: they leave equations unchanged, and a compound
formula stays compound.

For an arbitrary initial retained clause $\cl{A}{y.B}$, rather than one
generated by a proof rule, use its native representation
$A\Rightarrow_*\forall_*y\,B$. The derived implication and universal
eliminations simulate U. This extends the ordinary reading beyond the
origin-annotated generated contexts used in the executable translators.

\begin{theorem}[Canonical companion]
The ordinary realization of $\Can$ is bi-interpretable and equiconsistent
with that of $\Raw$. Its numerical interpretations use the whole domain
and actual equality. The raw-to-canonical numerical term composite is
provably equal to the identity; on canonical syntax the other composite
is the identity.
\end{theorem}
\begin{proof}
Compose the inherited zero-normalization equivalence, the canonical
equality bootstrap, and the clause/$Q$ equivalence just described.
The integration replaces the formula-context interface by clauses, retains
UE, and uses the inherited fixed-zero-base introduction with the complete
progress recovery just proved. All equality bootstrap proofs work under
arbitrary retained clauses. Formula transport and binder congruence extend
the numerical normalization identity through formulas.
\end{proof}


\section{Full iPRA and both interpretation composites}
\label{sec:interpretations}

\subsection{The ordinary theory presented}

Let $T_I$ be intuitionistic first-order logic with actual equality and
genuine falsity. Its function symbols are $0$, $S$, and the staged scalar
iterators: for each fixed earlier finite term $h(u,\vec p)$, add
$i_h(n,a,\vec p)$ with
\[
 i_h(0,a,\vec p)=a,\qquad
 i_h(Sn,a,\vec p)=h(i_h(n,a,\vec p),\vec p).
\]
Its remaining nonlogical principles are numerical coverage,
separation, and contextual induction on literal equations:
\[
 \forall x\,(x=0\lor\exists y\,x=Sy),\qquad
 (0=S0)\to\bot,\qquad\text{EI}.
\]
The base logic can be the usual intuitionistic natural deduction rules
for $\bot,\wedge,\vee,\to,\forall,\exists$ with reflexivity and equality
substitution. Its native realization is completely specified by
Section~\ref{sec:raw-spec}; no unspecified ordinary logic is required by
the standalone checker.

The ordinary readings of Eq, $D$, and clauses were given above. Under them
every rule is sound in $T_I$: Case uses coverage; the progressive branch
uses equation induction under the conditional base. Conversely, the
explicit logical macros give all intuitionistic rules, EI gives equation
induction, and numerical cases give coverage. Separation translates to
the identity implication $(0=S0)\Rightarrow_*(0=S0)$.

The falsity bridge must be explicit. Computation and positive equality
alone admit a singleton model; they cannot silently be identified with
a theory whose logical falsity has its usual interpretation.

\subsection{The inherited arithmetic bootstrap}

The included arithmetic proof proceeds in the following order:
\begin{enumerate}
\item Construct a scalar term $P$ with $P0=0$ and $P(Sn)=n$.
      This also proves successor injectivity.
\item Construct truncated subtraction and a faithful term test for
      actual numerical equality, valid at arbitrary model elements.
\item Turn each fixed quantifier-free formula into a characteristic
      equation. Its equation induction gives quantifier-free induction.
\item Construct pairing and projections by bounded search, and compile
      every fixed general primitive recursion into scalar iteration.
\end{enumerate}
The order matters: pairing and equality reflection are not assumed when
predecessor is constructed. The full inherited arguments appear in
Appendices~\ref{app:pred}--\ref{app:pr}. They are not claimed as new here.

For clarity, let $\pi,L,H$ be the constructed pairing and projections.
A target primitive recursion
\[
 f(0,\vec p)=g(\vec p),\qquad
 f(Sn,\vec p)=h(n,f(n,\vec p),\vec p)
\]
is compiled by
\[
\begin{aligned}
 V(n,\vec p)&=I\bigl(n,\pi(0,g^*(\vec p));
    u.\pi(SL(u),h^*(L(u),H(u),\vec p))\bigr),\\
 f^*(n,\vec p)&=H(V(n,\vec p)).
\end{aligned}
\]
Equation induction proves $L(V(n,\vec p))=n$, and the projection
equations prove the desired defining equations. This is an external
effective operation on fixed finite descriptions, not an internal
universal primitive-recursive evaluator.

\subsection{The composites are internally the identity}

Interpret an iterator in ordinary iPRA by primitive recursion whose
update ignores its counter. Interpret each fixed iPRA function description
in $T_I$ by the preceding compiler. In both directions the domain formula
is $x=x$, and equality remains $x=y$.

Equation induction proves uniqueness for each fixed scalar iteration and
each fixed ordinary primitive recursion. External induction on a finite
description shows that its double translation satisfies its own original
defining equations, hence is internally equal to the original function.
The graph $x=y$ is therefore an internally provable total bijection
preserving every operation in either numerical composite. This is a
statement about all numerical elements, not just standard numerals.

For formulas, the crucial composite is
\[
 \D_{x,y}B\quad\dashv\vdash\quad
 \exists_*x\bigl(B(x,0)\Rightarrow_*\forall_*y\,B(x,y)\bigr).
\tag{DC}\label{eq:dc}
\]
Forward, eliminate the paired formula at a fresh $b$, assume its base,
and use the resulting clause at an arbitrary numerical parameter; then
introduce the universal, implication, and existential macros. Backward,
eliminate the existential and introduce $D$ by Case, using the supplied
conditional universal at the successor. The ordinary interpretations of
the macros in \eqref{eq:logic} have their usual round trips by the
selector and gate endpoint proofs. Structural induction on formulas
completes both formula composites.

\begin{theorem}[Required equivalence]
The ordinary realizations of $\Raw$ and $\Can$ are each bi-interpretable
and equiconsistent with full first-order $\ipra$. Both numerical domains
are whole; interpreted equality is actual equality; both numerical
composite isomorphisms have graph $x=y$.
\end{theorem}
\begin{proof}
The preceding internal bridge gives $T_I\simeq\ipra$.
The inherited native logical bridge gives $\Old\simeq T_I$.
Theorem~\ref{thm:raw-equivalence} identifies every $\Raw$ sequent with
the corresponding $\Old$ sequent on unchanged syntax. The canonical
companion theorem provides the remaining equivalence. The explicit
falsity correspondence yields equiconsistency in both directions.
\end{proof}

\section{Complete checkers and measured costs}
\label{sec:implementation}

Each selected checker is one standalone Python module using only
\file{json} and \file{sys}. Its public input is exactly
\begin{lstlisting}
{"goal": <closed Eq/D formula>, "proof": <finite proof tree>}
\end{lstlisting}
Acceptance prints \file{ACCEPT} and exits with status zero. Rejection
prints \file{REJECT} and exits nonzero. Checks do not rely on Python
\file{assert}; optimized execution has the same intended acceptance
relation. The abstract calculus concerns arbitrary finite derivations;
ordinary finite-machine memory and recursion limits apply to the programs.

\subsection{Concrete representation and scope invariants}

Zero is integer $0$. Free eigenparameters are negative levels, and bound
variables use positive de Bruijn indices. Successor is \file{[0,t]},
iteration is \file{[n,a,h]}, equality is \file{[s,t]}, and a paired formula
is \file{[B]}. The iterator body binds one index; a paired formula binds
two. An opening replaces index $d+1$ at binder depth $d$ and decreases
larger indices by one. Its replacement is a well-scoped term with no
unbound positive indices, so negative free levels avoid capture without
shifting external parameters.

Terms, formulas, rule tags, rule arities, context indices, witnesses,
conclusions, and eigenparameter scopes are checked. Boolean JSON values
are not accepted as integers. Validation concerns the decoded JSON object;
the standard parser uses last-key semantics for duplicate textual keys. Each derivation rule preserves the typing
invariant of its inferred formula. The evaluator can therefore inspect
the shape of an already inferred equality without repeating a full
validator. A purported successor counter of any other shape is rejected.
The normal companion additionally rejects zero-counter iterators anywhere
in supplied term syntax and contracts them whenever opening creates one.

The finite CLI takes closed theorems from the empty context. The recursive
kernel implements local open sequents and clauses, which are used in every
introduction and elimination. The mathematical equivalence theorem is for
all open sequents, not merely the finite test inputs.

\subsection{Matching measurements}

Measurements use CPython 3.12.14. Remove the module docstring, parse the
entire source, print it with \file{ast.unparse}, and append one line feed.
Significant tokens exclude indentation, line-ending, end-marker, encoding,
and comment tokens. AST nodes count the entire parsed module. No imported
project helper or validator is excluded. Gzip uses level 9 and timestamp
zero, and is informational only.

\begin{center}
\small
\begin{tabular}{lrrrr}
\toprule
Complete checker & Bytes & Tokens & AST & Gzip\\
\midrule
Exact incoming raw clausal & 3046 & 853 & 993 & 990\\
Matched raw control & 3003 & 842 & 984 & 978\\
\textbf{Raw root-free proposal} & \textbf{2956} & \textbf{837} & \textbf{977} & \textbf{962}\\
\midrule
Zero-normal control with root case & 2993 & 864 & 1006 & 985\\
Zero-normal UE, witness guard & 2948 & 853 & 992 & 971\\
\textbf{Zero-normal UE, fixed-zero guard} & \textbf{2948} & \textbf{852} & \textbf{987} & \textbf{967}\\
Tagged native normal comparator & 3011 & 979 & 1047 & 965\\
\bottomrule
\end{tabular}
\end{center}

The matched raw control merely inlines a successor-constructor helper
that was called once. That cleanup is shared with the proposal.
Deleting the direct arithmetic branch then saves exactly 47 normalized
bytes (1.565\%), 5 significant tokens, and 7 AST nodes. Relative to the
unchanged incoming file the total savings are 90 bytes (2.955\%), 16
tokens, and 16 AST nodes; 43 of those bytes are the helper cleanup.

The witness-guard canonical intermediate saves 45 bytes, 11 tokens, and
14 AST nodes against its matching root-retaining control. Replacing its
witness guard by the fixed-zero-base guard then saves one token and five
AST nodes at unchanged byte count. The selected fixed-zero companion is
8 bytes smaller than the raw proposal but has 15 more tokens and 10 more
AST nodes. Consequently neither selected checker dominates the other on
all three coordinates.

The earlier tagged raw two-case checker has 2,996 bytes, 954 tokens,
and 1,019 AST nodes. The present raw clausal checker is smaller in all
three, but this comparison changes the logical interface and variable
representation; it is not the isolated 47-byte arithmetic ablation.
The repository's historical 516-byte compact profile and smaller validated
compressed encodings are not beaten. We claim a transparent complete
checker improvement, not a global source-size frontier.

\section{Verification, reproducibility, and limits}
\label{sec:verification}

The raw proposal was checked on 33 positive and 27 negative certificates,
including 120 ordinary and optimized command-line runs. All 33 positive
proofs passed the exact-sequent compiler, the reverse inclusion, and the
$N+2k$ node-count check. The examples include used induction hypotheses,
closure specialization, repeated-variable congruence, zero-source
transport, capture-avoiding witnesses, and nested compound binders.

The witness-guard canonical intermediate was checked on 15 positive and 27 negative
certificates, with 84 ordinary and optimized command-line runs. Each
positive proof was translated into the independently implemented inherited
tagged six-rule canonical checker and accepted. An additional 6,000
opening comparisons checked the negative-level representation against
that independent de Bruijn implementation. They include contractions
created by opening as well as binder-capture cases.

The selected fixed-zero companion was separately checked on 17 positive
and 28 negative certificates, with 90 ordinary and optimized command-line
runs. All 17 positive proofs were translated into the inherited seven-rule
zero-normal checker, which uses separate computation and transport rules,
and accepted. Additional cases use progress at a nonzero witness, cases
at witness zero with base $1=1$, and rejection of progress when the base
is not $0=0$. Thus the changed guard is tested directly.

Invalid cases include escaping eigenparameters, an unshifted or unbound
index, invalid iterator counters, forged evaluation sources, a wrong
clause antecedent, Boolean indices, compound-progress attempts, malformed
proof arities, and extra top-level fields. There were zero unexpected
outcomes in the recorded runs.

These are implementation checks. The universal correctness statements
are the written derivation transformations and inherited internal
arithmetic proofs. This package does not claim a new Lean or Rocq proof.
The upstream README reports a compiled Lean development, but it was not
rebuilt during this work. The selected raw and canonical Python profiles
have different syntactic contracts; no general raw-to-canonical proof
compiler is advertised. The executable raw translation stays within raw
syntax, and the canonical translation targets the separate inherited
canonical checker.

From the extracted package root, run
\begin{lstlisting}
python3 reproduce.py
python3 code/raw_two_case.py < code/raw_certificates/used_induction_hypothesis.json
python3 code/zero_base.py < code/zero_base_certificates/used_equation_induction.json
pdflatex -interaction=nonstopmode -halt-on-error report.tex
pdflatex -interaction=nonstopmode -halt-on-error report.tex
\end{lstlisting}
The report's distributed \file{report.tex} is standalone: its inherited
appendices and code listings are embedded. The editable assembly source
and source fragments are also provided. The checks require only Python's
standard library; the PDF rebuild requires a normal TeX Live installation.

\section{Further research questions}
\label{sec:questions}

\begin{enumerate}
\item \textbf{Formalize the exact integration.} Add a clausal root-free
      calculus to the existing Lean development and prove the sequent
      equivalence by the three-node frame replacement. This is a small
      extension relative to the arithmetic bridge, but the context clauses
      and their occurrence annotations should be represented explicitly.
\item \textbf{Can the clean rule basis improve the compact frontier?}
      The present matched strict improvement does not beat the old compressed
      codecs. A meaningful further reduction must count grammar validation,
      normalization if used, and the complete public interface.
\item \textbf{Can introduction be made both simpler and substitution-stable?}
      The raw profile retains an explicit literal-witness test; the canonical
      profile selects progress by the exact zero base. The existing diagonal
      guard is stable in both directions under substitution, while the
      fixed-zero guard may change from false to true. Can this remaining
      proof-translation complication be removed at equal or smaller cost?
\item \textbf{Can proof certificates share the iterator frames?}
      The $N+2k$ theorem controls tree nodes. DAG certificates could share
      repeated frames and reflexivities, but their reference validation and
      acyclicity checks could erase the verifier saving.
\item \textbf{Can a narrower uniqueness rule replace equation induction?}
      The incoming scalar-uniqueness package establishes a full scalar
      uniqueness bridge. The separate successor-rigidity report reaches
      zero-equation induction and leaves full equality reflection open.
      Neither fact licenses removing the equation-progress restriction here.
\item \textbf{What cost model should define a genuine minimum?}
      Rule rows, trusted normalization, parser scope, certificate format,
      source bytes, AST nodes, and running time measure different things.
      A formal optimization problem needs to fix these choices before an
      absolute simplest-presentation claim can have a determinate meaning.
\end{enumerate}

\section*{Attribution and references}
\addcontentsline{toc}{section}{Attribution and references}
The clause interface, scalar bridge, zero normalization, witness-directed
introduction, and iterator-frame equality reductions are inherited from
ProveIt contributors. Code and reproduced mathematical source fragments
retain the accompanying MIT No Attribution license. The new work is the
combined specification, matching implementation, proof compilers,
comparison, and validation described above.

\begin{thebibliography}{9}
\small\raggedright
\bibitem{baseline} ProveIt contributors. \emph{Concise presentations of
intuitionistic primitive recursive arithmetic}. Pinned snapshot
\texttt{831081033253cca2115196cf86bb99c77ed2aeae}.
\url{https://github.com/VladimirReshetnikov/ProveIt/tree/831081033253cca2115196cf86bb99c77ed2aeae/Logic/PrimitiveRecursiveArithmetic/ConciseIPRA}.
\bibitem{raw} ProveIt contributors. \emph{Five rules for iPRA with ordinary
substitution}. \file{docs/incoming/ipra_raw_clausal_2026-11-11.zip}, same snapshot.
\bibitem{iteration} ProveIt contributors. \emph{Equality from iteration}.
\file{docs/incoming/ipra_equality_from_iteration_2026-11-11.zip}, same snapshot.
\bibitem{canonical} ProveIt contributors. \emph{Five rules for intuitionistic
primitive recursive arithmetic}. \file{docs/incoming/ipra_clausal_canonical.zip}, same snapshot.
\bibitem{reflexive} ProveIt contributors. \emph{Reflexive base}.
\file{docs/incoming/ipra_reflexive_base.zip}, same snapshot.
\bibitem{unique} ProveIt contributors. \emph{Scalar iteration uniqueness for
iPRA}. \file{docs/incoming/ipra_scalar_uniqueness.zip}, same snapshot.
\bibitem{native} ProveIt contributors. \file{Lean/ConciseIPRA/NativeEval.lean},
\file{NativeZeroNormal.lean}, \file{NativeBiInterpretation.lean},
and \file{BiInterpretation.lean}, same project snapshot.
\bibitem{bridge} ProveIt contributors. Five internal arithmetic and logical
bridge fragments under \file{Research/certified_zero/reference/},
reproduced in the appendices with their license.
\bibitem{erasure} ProveIt contributors. \emph{Internal scalar predecessor
bootstrap}. \file{Research/certified_zero/reference/scalar_predecessor.tex},
same snapshot; reproduced here.
\bibitem{forward} ProveIt contributors. \emph{From literal-equation induction
to full intuitionistic PRA}.
\file{Research/certified_zero/reference/equality_reflection.tex},
same snapshot; reproduced here.
\bibitem{progress} ProveIt contributors. \emph{Pairing, primitive recursion,
and internal numerical composites}.
\file{Research/certified_zero/reference/primitive_recursion.tex},
same snapshot; reproduced here.
\end{thebibliography}

\clearpage
\appendix
\part*{Inherited bridge proofs}
\addcontentsline{toc}{part}{Inherited bridge proofs}
The following five source fragments are inherited reference material.
Historical words such as ``new'' and ``old'' inside a fragment refer to
that fragment's original comparison, not to novelty claimed by this
report. Its clause basis has both explicit Case and equation-progress
choices; Section~\ref{sec:recovery} explains their recovery in the selected
five-pattern kernel. The entire dependency chain is included so the
bi-interpretation claim does not rest on agreement over standard numerals.

% New mathematical result: complementary introduction choices.
% Requires amsmath/amssymb/amsthm and theorem, lemma, proposition environments.
% The inherited arithmetic-to-iPRA bridge must be supplied separately.

\section{Complementary introductions with one local assumption}
\label{app:choice}
\label{sec:choice-introduction}

The old paired introduction retains both the zero instance and an
induction instance when its matrix is an equation. Both together are
unnecessary. A certificate may choose either one, provided the induction
choice is restricted to a literal equation. This changes neither the
formula language nor the meaning of a formula. The omitted combination
is recovered by a short nested introduction--elimination construction.

\subsection{Complete native syntax and rule inventory}

The scalar numerical language is
\[
 t::=x\mid0\mid St\mid I(n,a;u.h).
\]
Only $h$ binds the state variable $u$; $n$ and $a$ do not.
The update may contain fixed numerical parameters. The computation
axioms in every retained context are
\[
 I(0,a;u.h)=a,\qquad
 I(Sn,a;u.h)=h[I(n,a;u.h)/u].
 \tag{I0, IS}\label{eq:choice-computations}
\]
The argument also applies unchanged to the full fixed-description PR
language with its projection, composition, and recursion equations.
No numerical variable ranges over function descriptions, and no uniform
evaluator is postulated.

Native formulas and reusable context clauses are
\[
 A::=(s=t)\mid\mathsf D_{x,y}A,\qquad
 \langle A;y.B(y)\rangle,\quad y\notin\operatorname{FV}(A).
\]
Bound variables are renamed before capture-free substitution.
The two variables $x,y$ in a paired formula are distinct.
An ordinary assumption $F$ is the constant clause
$\langle H;v.F\rangle$, where $v$ is absent from $F$ and $H$ is the
fixed computation theorem $I(0,0;u.0)=0$. Thus ordinary hypothesis
use is a clause Use preceded by this computation proof.

There are computation axioms and exactly the following inference
patterns:
\[
\frac{\Gamma\vdash A}{\Gamma\vdash B(t)}
\quad\mathsf{Use}
\quad(\langle A;y.B(y)\rangle\in\Gamma),
\]
\[
\frac{\Gamma\vdash\mathsf D_{x,y}B(x,y)\qquad
      \Gamma,\langle B(a,0);y.B(a,y)\rangle\vdash F}
     {\Gamma\vdash F}
\quad\mathsf{DE},
\qquad
\frac{\Gamma\vdash s=t\qquad\Gamma\vdash F(s)}
     {\Gamma\vdash F(t)}
\quad\mathsf{Tr},
\]
\[
\frac{\Gamma,B(t,\eta_\varepsilon(z))\vdash B(t,Sz)}
     {\Gamma\vdash\mathsf D_{x,y}B(x,y)}
\quad\mathsf{DI}_{\varepsilon},
\qquad
\eta_0(z)=0,\quad\eta_1(z)=z.
\tag{Choice DI}\label{eq:choice-di}
\]
Here $\varepsilon$ is a certificate choice in $\{0,1\}$, and
$\varepsilon=1$ is permitted only when $B$ has a literal equation
root. An equation beneath a binder is not a literal equation.
Call $\varepsilon=0$ the \emph{case choice} and $\varepsilon=1$
the \emph{progress choice}. Both retain exactly one ordinary
formula assumption. Compound matrices have only the case choice.

In DI the eigenparameter $z$ is fresh for $\Gamma$, the witness
$t$, and every fixed parameter of $B$. In DE the eigenparameter
$a$ is fresh for $\Gamma$, the conclusion $F$, and every fixed
parameter of $B$. In transport, $F(v)$ is any fixed native
formula template. Displayed local assumptions and clauses are
discharged. There are no primitive reflexivity, coverage,
separation, ordinary logical, or additional induction rules.

We call this calculus $\mathcal C$. The original comparison
calculus $\mathcal N$ has identical syntax, computation,
Use, DE, and transport, but its introduction is
\[
\frac{\Gamma,B(t,0),[B(t,z)]_{\mathrm{at}}\vdash B(t,Sz)}
     {\Gamma\vdash\mathsf D_{x,y}B(x,y)}.
\tag{Old DI}\label{eq:choice-old-di}
\]
The bracket supplies $B(t,z)$ exactly at a literal equation.

\subsection{Exact equivalence and the unused-witness construction}

\begin{lemma}[Structural admissibility]
\label{lem:choice-structural}
Weakening, exchange, contraction, renaming, ordinary-assumption cut,
and capture-free numerical substitution are admissible in $\mathcal C$.
\end{lemma}
\begin{proof}
Induct on derivations. Rename local eigenparameters before inserting
proofs or terms. Numerical substitution preserves a formula's root
constructor, so an allowed progress choice remains allowed.
The selected choice does not change under substitution or cut.
For ordinary-assumption cut, replace each use of its constant clause
by the supplied proof. Computation schemas are closed under
capture-free numerical substitution. These facts verify every case.
\end{proof}

\begin{theorem}[Identity-language equivalence]
\label{thm:choice-equivalence}
For every retained clause context $\Gamma$ and native formula $F$,
\[
 \Gamma\vdash_{\mathcal C}F
 \quad\Longleftrightarrow\quad
 \Gamma\vdash_{\mathcal N}F.
\]
The transformations are effective on finite derivations and leave
all numerical and formula syntax unchanged.
\end{theorem}
\begin{proof}
Each new introduction is an instance of the old one after weakening
by the omitted local assumption. For the case choice at an equation,
add the induction instance. For the progress choice, add the zero
instance. Compound case introduction is already the old rule.
Thus every new derivation translates to an old one.

For the reverse direction, transform an old derivation recursively.
Only its atomic DI case needs a construction. Suppose its matrix is
the equation $E(x,y)$ and its transformed premise is
\[
 \pi:\quad\Gamma,E(t,0),E(t,z)\vdash E(t,Sz).
\]
The outer new inference will be case DI with witness $t$. It is
therefore enough to prove $E(t,Sz)$ under
$\Gamma,E(t,0)$.

Rename the eigenvariable in $\pi$ to a fresh $v$. Choose a new
variable $u$ absent from $E(t,y)$ and all retained data. Under
$\Gamma,E(t,0)$, progress DI with witness zero introduces
\[
             \mathsf D_{u,y}E(t,y)
\]
from exactly the renamed premise $\pi[v/z]$.
The bound coordinate $u$ is deliberately unused.

Apply DE to this paired formula. Because $u$ is unused, the opened
clause is exactly
\[
               \langle E(t,0);y.E(t,y)\rangle;
\]
it contains no new existential witness. The retained ordinary
assumption $E(t,0)$ supplies its antecedent. Use at $Sz$ gives
$E(t,Sz)$, and DE closes. Finally the outer case DI introduces
the original $\mathsf D_{x,y}E(x,y)$ with witness $t$.
The original premise proof $\pi$ has been inserted once.

An old compound introduction already is the new case choice.
All other rule cases are unchanged. This completes the induction
on the original derivation.
\end{proof}

The unused inner coordinate is essential. Introducing
$\mathsf D_{x,y}E(x,y)$ internally and then eliminating it would
give an arbitrary fresh witness, not necessarily the original
$t$. The construction instead keeps $t$ as a fixed parameter
of $E(t,y)$ and binds only the unused coordinate $u$.

\begin{proposition}[Linear bound on inference nodes]
\label{prop:choice-node-bound}
Count each computation axiom, Use, DI, DE, and transport as one
inference node, expanding ordinary-assumption use as a computation
of $H$ followed by Use. If an old derivation has $N$ nodes and
$r$ atomic DI nodes, the displayed transformation has exactly
\[
                   N+5r\leq 6N
\]
nodes.
\end{proposition}
\begin{proof}
Each atomic old DI node is replaced by six nodes: outer case DI,
inner progress DI, DE, Use of the opened clause, Use of the
retained base assumption, and the computation of $H$ needed for
that last ordinary-assumption use. Its old premise subtree
occurs once. All other inference nodes are copied once.
\end{proof}

This counts inference-tree nodes. It is not a bound on serialized
JSON bytes, formula sizes, checker runtime, or the number of
consumed packet cells when hypothesis closure is implicit.
In the formula-context Q/D presentation, the corresponding
replacement has outer case DI, inner progress DI, DE, QE,
and two Hyp nodes, giving the same six-node count.

\subsection{Induction and numerical cases separate cleanly}

\begin{lemma}[Contextual equation induction]
\label{lem:choice-ei}
The progress choice, DE, and Use derive
\[
\frac{\Gamma\vdash E(0)\qquad
      \Gamma,E(z)\vdash E(Sz)}
     {\Gamma\vdash E(t)}
\tag{EI}\label{eq:choice-ei}
\]
for every equation $E$, with the usual freshness condition.
\end{lemma}
\begin{proof}
Choose an unused witness coordinate $u$. Progress DI introduces
$\mathsf D_{u,y}E(y)$ from the given step. DE opens the clause
$\langle E(0);y.E(y)\rangle$. Supply its antecedent using the given
base proof and use the clause at $t$. Close DE.
\end{proof}

\begin{lemma}[Cases, equality, and binder congruence]
\label{lem:choice-preinduction}
The case choice derives, for arbitrary native formulas,
\[
\frac{\Gamma\vdash F(0)\qquad\Gamma\vdash F(Sz)}
     {\Gamma\vdash F(t)}
\tag{NC}\label{eq:choice-nc}
\]
with the usual freshness condition. Equality reflexivity,
symmetry, transitivity, numerical congruence, and uniform
matrix interderivability under $\mathsf D$ are derivable
without using the progress choice.
\end{lemma}
\begin{proof}
For cases, use case DI with an unused witness coordinate and
matrix $F(y)$, ignoring its retained base in the successor
premise. DE opens $\langle F(0);y.F(y)\rangle$; the given zero
proof and Use at $t$ finish.

For reflexivity put $r=I(0,t;u.0)$. Use the computation equation
$r=t$ as both transport premises with motive $v=t$ to obtain
$t=t$. The usual explicit equality motives give symmetry,
transitivity, and numerical congruence.

For binder congruence, open the old paired formula at a fresh
witness $a$ and introduce the new one by case DI at $a$.
Convert its retained base to the old base, use the opened
old clause at the successor, and convert back. Ordinary cut
closes the construction, and the reverse conversion gives
the converse. No induction is used.
\end{proof}

The primitive progress choice is still restricted to equations.
Allowing it at arbitrary compound matrices would, by the
argument for \eqref{eq:choice-ei}, supply induction for
arbitrary formulas and change the intended theory.

\subsection{The entire ordinary logical interface}
\label{subsec:choice-logic}

Here is a predecessor-free reconstruction of the ordinary logical
operations. It makes explicit that the case choice supplies the
logical interface independently of progress induction.

Let $C(v,a,b):=I(v,a;u.b)$ with $u$ unused. Computation gives
$C(0,a,b)=a$ and $C(Sv,a,b)=b$. Define the gate on fixed finite
formula syntax by
\[
 G_v(s=t):=(C(v,0,s)=C(v,0,t)),\qquad
 G_v(\mathsf D_{x,y}B):=\mathsf D_{x,y}G_vB,
\]
renaming bound variables away from the control. Equality,
binder congruence, and external structural induction give
\[
 \vdash G_0A,\qquad G_{Sv}A\dashv\vdash A,\qquad
 A\dashv\vdash\mathsf D_{u,y}G_yA
 \quad(u,y\notin\operatorname{FV}(A)).
\tag{Gate and padding}\label{eq:choice-gate}
\]
For padding, case DI proves the forward direction. DE, the
zero-gate theorem, and Use at $1$ prove the reverse direction.

Define a selector between two arbitrary formulas by
\[
 M_v(s=t,u=w):=(C(v,s,u)=C(v,t,w))
\]
and, with aligned fresh paired variables, by
\[
\begin{aligned}
 M_v(\mathsf D_{x,y}A,\mathsf D_{x,y}B)
   &:=\mathsf D_{x,y}M_v(A,B),\\
 M_v(E,\mathsf D_{x,y}B)
   &:=\mathsf D_{x,y}M_v(G_yE,B),\\
 M_v(\mathsf D_{x,y}A,E)
   &:=\mathsf D_{x,y}M_v(A,G_yE),
\end{aligned}
\]
where $E$ is an equation. Every call decreases maximum input
binder depth. Padding and congruence prove
$M_0(A,B)\dashv\vdash A$ and
$M_{Sv}(A,B)\dashv\vdash B$.

Put $\delta(v)=C(v,1,0)$ and $p(z)=I(z,0;u.\delta(u))$.
Only the fixed computations $p(1)=1$ and $p(2)=0$ are needed.
The full logical macro inventory is
\[
\begin{aligned}
 A\Rightarrow_*B&:=\mathsf D_{u,z}M_z(A,B),\\
 A\wedge_*B&:=\mathsf D_{u,z}G_zM_{p(z)}(A,B),\\
 \exists_*x\,A(x)&:=\mathsf D_{x,y}G_yA(x),\\
 A\vee_*B&:=\exists_*x\,M_x(A,B),\\
 \forall_*x\,A(x)&:=
        A(0)\wedge_*\mathsf D_{u,z}G_zA(z),\\
 \bot_*&:=(0=1).
\end{aligned}
\tag{Logical macros}\label{eq:choice-macros}
\]
Newly bound variables are fresh; $u$ is unused.

Implication introduction obtains its assumption from the zero
selector and supplies the consequent at successors, then uses
case DI. Elimination opens its clause and uses it at $1$.
Existential introduction uses its witness and the successor
gate; elimination opens at a fresh witness, supplies the
zero-gate theorem, and uses the clause at $1$.

Disjunction introduction uses witnesses $0$ and $1$.
For elimination, discharge the two branch proofs as
$A\Rightarrow_*F$ and $B\Rightarrow_*F$. NC applied to
$M_z(A,B)\Rightarrow_*F$ gives this implication at every
control: the two endpoint cases follow from the branch
implications. Existential elimination provides $M_a(A,B)$
and the implication gives $F$.

Given $A$ and $B$, selector endpoints and NC give $M_v(A,B)$
at an arbitrary control. Substitute $v=p(Sz)$ and use the
successor gate for conjunction introduction. For projections,
open its zero-gated clause and specialize at $1$ and $2$;
the two fixed computations give $B$ and $A$, respectively.

These rules prove existential numerical coverage:
\[
     (x=0)\vee_*\exists_*v\,(x=Sv).
\tag{Coverage}\label{eq:choice-coverage}
\]
Use NC. Zero takes the left injection of reflexivity, and a
successor takes the right injection with the evident witness.
Coverage gives equality cases even when the retained context
contains $x$; its elimination adds $x=0$ or $x=Sv$ rather
than silently changing that retained context.

Universal introduction follows from a fresh-parameter proof
by substituting zero and successors and forming the two
conjuncts. For elimination, take its $A(0)$ conjunct and
open the paired conjunct. Its base is a zero-gate theorem.
Coverage of the target chooses the zero instance from
$A(0)$ or a successor instance from the opened clause
and gate removal.

Finally, under $0=1$, computation and transport give
$s=C(0,s,t)=C(1,s,t)=t$ for arbitrary numerical terms.
Every equation follows, and external induction on formula
syntax with case DI gives every native formula. This
supplies falsity elimination. No excluded middle for
quantified formulas has entered the construction.

\subsection{Whole-domain realizations and both formula composites}

For a standard first-order reading, let $T_{\rm eq}$ have
intuitionistic equality logic, the same numerical signature
and computation equations, numerical coverage, separation
$(0=1)\to\bot$, and contextual equation induction.
Its interpretation of the native syntax is
\[
 (s=t)^\circ=(s=t),\qquad
 (\mathsf D_{x,y}B)^\circ
   =\exists x\bigl(B^\circ(x,0)\to\forall y B^\circ(x,y)\bigr),
\]
\[
 \langle A;y.B\rangle^\circ
       =A^\circ\to\forall y B^\circ.
\]
The domain is the entire numerical domain and interpreted
equality is actual equality.

Case DI is sound in this realization by numerical coverage.
Progress DI is sound by contextual equation induction under
the conditional base. Use, DE, transport, and computation
have their displayed ordinary meanings. Conversely, the
native macros derive every ordinary logical rule,
\eqref{eq:choice-ei} derives equation induction, and
\eqref{eq:choice-coverage} derives coverage. Separation
translates to the identity implication
$(0=1)\Rightarrow_*(0=1)$.

For the essential paired formula composite,
\[
 \mathsf D_{x,y}B
 \dashv\vdash
 \exists_*x\bigl(B(x,0)\Rightarrow_*\forall_*y\,B(x,y)\bigr).
\tag{Paired composite}\label{eq:choice-paired-composite}
\]
Forwards, DE opens its clause at a fresh $a$; under $B(a,0)$
that clause supplies $B(a,y)$ at a fresh arbitrary $y$.
Introduce the universal, implication, and existential
macros and close DE. Backwards, eliminate the existential
at a fresh $a$. Introduce the paired formula with witness
$a$ by case DI: the retained base, supplied implication,
and universal elimination at $Sz$ give its successor premise.

The ordinary interpretations of the six logical macros
have their ordinary round trips. For conjunction, its true
base and the instances $1,2$ recover both operands;
the converse uses selector cases. The universal has all
successor instances in its paired conjunct and zero in
its first conjunct; coverage gives exactly the ordinary
universal. Implication, disjunction, and existential
quantification use their selector or gate endpoints.
Falsity uses the explicit ordinary separation bridge.
External induction on finite formulas, together with
\eqref{eq:choice-paired-composite}, gives both formula
composites. Clauses correspond in both directions to
$A\Rightarrow_*\forall_*yB$ by the derived logical rules.

The numerical composites are literal identities with
internally provable graph $x=y$. No standard-numeral
restriction, retract, quotient, or equality reflection
assumption is being inserted at this stage. The separately
supplied arithmetic bridge identifies $T_{\rm eq}$ with
full iPRA; the present change preserves that bridge
unchanged because Theorem~\ref{thm:choice-equivalence}
preserves every native sequent.

\subsection{What the theorem does and does not claim}

The new inference interface retains one ordinary assumption
at every paired introduction. It preserves the entire
consequence relation, not just equations or closed
theorems. Consequently the established full-iPRA
interpretations and equiconsistency transfer unchanged.

This result is different from retaining only progress
introduction and strengthening closure elimination.
Here the original elimination rule is retained, and
the case and progress introductions coexist. It is also
different from allowing arbitrary compound progress,
which would supply stronger induction.

A complete-checker comparison must still count the
certificate choice, its equation-root restriction,
substitution, context handling, computation, and the
public input/output interface. The mathematical theorem
alone is not a byte-count or minimality theorem.
The inference-node bound is distinct from source-size
measurements and from serialized certificate size.
The proof is a written universal proof; finite regression
tests do not constitute a Lean/Rocq formalization of it.
\section{Arithmetic strength and both ordinary interpretation composites}
\label{app:ordinary}
\label{sec:ordinary-bridge}

We spell out the arithmetic bridge to avoid identifying standard numerical
agreement with bi-interpretability. First elaborate every fixed iterator
into an ordinary function symbol. At stage zero the symbols are $0,S$.
Given a fixed earlier-stage term $h(u,\vec p)$, introduce
$i_h(n,a,\vec p)$ with the two equations
\[
i_h(0,a,\vec p)=a,\qquad
i_h(Sn,a,\vec p)=h(i_h(n,a,\vec p),\vec p).
\]
Only finite descriptions occur; the domain contains no code for a
universal evaluator. Formal parameters are distinct when a description
is declared, while actual arguments can coincide.

Let $T_I$ be ordinary intuitionistic first-order logic with actual
equality and genuine falsity, these staged equations, contextual
induction for literal equations, and
\[
\forall x\,(x=0\vee\exists y\,x=Sy),
\qquad (0=S0)\to\bot.
\]
The second sentence is the falsity bridge. It must be stated: positive
numerical equations by themselves admit a singleton structure. Native
explosion from $0=S0$ and a translation of ordinary falsity to that
equation do not authorize silently dropping the bridge from $T_I$.

The native-to-ordinary translation is
\[
(s=t)^\circ=(s=t),\qquad
(\mathsf D_{x,y}B)^\circ=
\exists x\bigl(B^\circ(x,0)\to\forall y B^\circ(x,y)\bigr).
\]
A clause $\langle A;y.B\rangle$ becomes $A^\circ\to\forall yB^\circ$.
For the auxiliary closure used by the compact checkers,
$(\mathsf Q_yB)^\circ=B^\circ(0)\to\forall yB^\circ$.
The interpretation domain is all elements, defined by $x=x$;
equality is actual equality.

Case DI is sound because numerical coverage combines its retained base
with its successor proof. Progress DI is sound by literal-equation
induction under a temporarily assumed base. Use, DE and transport are
ordinary logical inferences under the translations. Conversely the
derived logical operations, equation induction, cases and explosion in
Section~\ref{sec:choice-introduction} recover the rules and axioms of
$T_I$. An auxiliary closure is interderivable with
$\mathsf D_{v,y}B(y)$ for unused $v$: in one direction use case DI with
the closure and its base; in the other eliminate the paired formula
and return its supplied closure. Thus formulas and context clauses have
the stated meanings in both presentations.

\subsection{The arithmetic dependency order}

The scalar update has no recursion-counter argument. We do not conceal
one in an encoding. The inherited internal proofs, reproduced in the
appendices, establish the required arithmetic in this order:
\begin{enumerate}
\item Using equation induction and numerical cases, construct a fixed
scalar predecessor term and prove $P0=0$ and $P(Sn)=n$.
Existential coverage then gives $x=0\vee x=S(Px)$: in the successor
case $x=Sy$, apply $P$ to obtain $Px=y$ and substitute.
\item Construct truncated subtraction and a numerical characteristic
of actual equality. Prove its reflection on the entire original domain.
\item Compile each fixed quantifier-free formula into its characteristic
equation and recover every contextual quantifier-free induction instance.
\item With that induction available, construct pairing and projections
by a bounded square-root search and compile every fixed general
primitive-recursive description into a scalar iterator term.
\end{enumerate}
No pairing or faithful equality detector is assumed to obtain
predecessor. No appeal to standard-model correctness replaces an
internal equality. Every induction fixes its passive parameters;
diagonal substitution occurs only into a previously proved identity.

For a target recursion
\[
f(0,\vec p)=g(\vec p),\qquad
f(Sn,\vec p)=h(n,f(n,\vec p),\vec p),
\]
let $\pi,L,H$ be the proved pairing and projections. The compiler uses
\[
V(n,\vec p)=I\bigl(n,\pi(0,g^*(\vec p));
u.\pi(SL(u),h^*(L(u),H(u),\vec p))\bigr),
\qquad f^*(n,\vec p)=H(V(n,\vec p)).
\]
Equation induction proves $L(V(n,\vec p))=n$; the two projection
identities then prove the target defining equations. The construction
is external recursion on each finite PR description, not an internal
quantification over arbitrary functions.

\begin{theorem}[Whole-domain bi-interpretability and equiconsistency]
The ordinary realization $T_I$ and full first-order iPRA are
bi-interpretable. In both theories $x=y$ defines the isomorphism from
the original domain to its double interpretation. The native Choice-DI
calculus is therefore a faithful proof presentation of a theory
bi-interpretable and equiconsistent with iPRA.
\end{theorem}
\begin{proof}
Interpret an iterator in iPRA by primitive recursion that ignores its
counter argument in the update. iPRA proves its equations, literal
equation induction, coverage and separation. In the reverse direction
use the scalar compiler. The predecessor and equality-reflection proofs
supply the successor axioms and every quantifier-free induction instance.
All ordinary logical rules are preserved.

For two iterators with the same seed and provably equal updates,
equation induction on the equality of their values proves uniqueness at
every counter, with all parameters fixed. The corresponding argument
applies to ordinary primitive recursion. Induct externally on a finite
description: its double translation satisfies its original defining
equations, so internal uniqueness identifies the two functions. This
also proves coherence when substitution changes the staged syntactic
description of an update.

The domain and equality are unchanged; hence the identity graph is
internally total, bijective and preserves every interpreted operation.
The displayed ordinary meanings of the logical macros and the paired
binder give both formula composites by external induction on finite
formula syntax. Clauses are similarly equivalent to their conditional
universal meanings. The explicit falsity bridge preserves contradiction
in both directions. Effective translations of finite proofs give
equiconsistency, without a claim that either theory proves its own
consistency.
\end{proof}
% Adapted with attribution from the pinned unparenthesized one-hole archive.
\section{Internal scalar predecessor bootstrap}
\label{app:pred}
\label{sec:hybrid-bootstrap}

This scalar bootstrap is inherited from the supplied scalar development~\cite{erasure} and is reproduced to keep the scalar checker self-contained. It is not a new contribution of this paper. The proof uses a factor map fixing odd states and a saturated-iteration construction of predecessor.

The available arithmetic terms are $0$, $S$, and
\[
 I(n,a;u.h),\qquad
 I(0,a;u.h)=a,\qquad
 I(Sn,a;u.h)=h[I(n,a;u.h)/u].
\]
The update binds only $u$.  All other variables in it are fixed numerical
parameters.  We use equality reasoning, the already reconstructed
zero/successor cases, and induction on a single literal equation.
Every induction below keeps its other parameters fixed.  Predecessor,
successor injectivity, pairing, and induction on arbitrary
quantifier-free formulas are not assumptions of the argument.

\subsection{Elementary terms and two finite cycles}

Write $1=S0$ and use larger numerals as successor abbreviations.  Define
\[
 C(n,a,b)=I(n,a;u.b),\qquad
 Z(n)=C(n,1,0),\qquad
 a+b=I(b,a;u.Su).
\]
The update variable of $C$ is absent from its body.  Consequently
\[
 C(0,a,b)=a,\qquad C(Sn,a,b)=b.
\]
Numerical cases show that $Z(n)$ is $0$ or $1$.
The usual commutative-monoid equations for addition follow by literal
equation induction.  For example, $0+b=b$ and $Sa+b=S(a+b)$ are
proved by induction on $b$; these yield commutativity by a further
equation induction, and associativity follows by induction on the last
summand.  We use these equations to reorder finite sums.  The notation
$2a$ means $a+a$ and $3a$ means $a+a+a$; no multiplication operation
is needed for this bootstrap.

Put
\[
 N(n)=I(n,0;u.Z(u)),\qquad
 j(u)=Z(u)+2N(u),\qquad
 M(n)=I(n,0;u.j(u)).
\]
The range of $N$ is contained in $\{0,1\}$: its zero value is $0$,
and its successor value is a value of $Z$.  This uses numerical cases
on its counter, not induction on a disjunction.  Hence
\begin{equation}
 N(n)+Z(N(n))=1,\qquad
 Z(Z(N(n)))=N(n),\qquad N(n+2)=N(n).
 \label{eq:hybrid-parity}
\end{equation}
Here the last equation follows from two successor computations and
the preceding equation.  Literal-equation inductions give
\begin{equation}
 N(2n)=0,\qquad N(2n+1)=1,\qquad N(3n)=N(n).
 \label{eq:hybrid-parity-scaling}
\end{equation}
For the first equation the successor step adds two to the argument of
$N$; the second follows by one computation.  In the third, increasing
$n$ adds three to the left argument and therefore applies $Z$ once to
both sides of the induction equation.

Finite computations now give
\[
 j(0)=1,\qquad j(1)=2,\qquad j(2)=0.
\]
The proof below needs just these three values and
$M(Sn)=j(M(n))$; it does not assume a division algorithm or a theorem
characterizing $M$ as remainder modulo three.

\subsection{A ramp with a global scaling identity}

Define
\[
 G(n)=I(n,0;u.u+M(u)+1).
\]
The central invariant is
\begin{equation}
 M(G(n))=N(n).
 \label{eq:hybrid-ramp-invariant}
\end{equation}
Its zero case is immediate.  For its successor step, retain the
single induction equation and split the Boolean value $N(n)$.
If it is $0$, the next $G$ state is $G(n)+1$; its $M$ value is
$j(0)=1$, equal to $N(Sn)$.  If it is $1$, the next $G$ state
is $G(n)+2$; two computations of $M$ change its value from $1$
to $2$ and then to $0$, again equal to $N(Sn)$.  This proves
the invariant by literal-equation induction.

It follows that
\[
 G(Sn)=G(n)+N(n)+1,\qquad G(n+2)=G(n)+3.
\]
Equation induction therefore yields
\begin{equation}
 G(2n)=3n.
 \label{eq:hybrid-ramp-double}
\end{equation}
A separate literal-equation induction proves
\begin{equation}
 G(3n)=3G(n)+N(n).
 \label{eq:hybrid-ramp-triple}
\end{equation}
For completeness, its successor calculation is
\begin{align*}
 G(3Sn)
 &=G(3n)+N(3n)+4
   =3G(n)+2N(n)+4,\\
 3G(Sn)+N(Sn)
 &=3G(n)+3N(n)+3+Z(N(n)).
\end{align*}
The expressions agree in both Boolean cases for $N(n)$.
The base is $0=0$.  In particular,
\eqref{eq:hybrid-ramp-triple} was proved internally; it was not
inferred from an intended floor-function interpretation of $G$.

Set
\[
 H(u)=C(N(u),G(u),u).
\]
The crucial choice is to retain $u$ when its parity is $1$.
The preceding equations give
\begin{equation}
 H(2u)=3u,\qquad H(1)=1,\qquad H(3u)=3H(u).
 \label{eq:hybrid-factor-map}
\end{equation}
To check the last identity, use $N(3u)=N(u)$ and split the Boolean
$N(u)$.  At $0$, both sides are $3G(u)$ by
\eqref{eq:hybrid-ramp-triple}; at $1$, both are $3u$.
This argument does not require that an arbitrary number has already
been represented as $2v$ or $2v+1$.

\subsection{A comparator from scalar iteration}

For every fixed update $h$, equation induction on $n$, keeping $m,a$
and all parameters fixed, proves the semigroup identity
\[
 I(m+n,a;h)=I(n,I(m,a;h);h).
\]
The base is the zero computation and the step applies $h$ to the
induction equation. This identity will also be used in the later
pairing construction.


For a fixed update $h$, equation induction on $y$, with its initializer
$a$ and all parameters of $h$ fixed, proves the rotation identity
\begin{equation}
 I(Sy,a;h)=I(y,h(a);h).
 \label{eq:hybrid-rotation}
\end{equation}
Both sides have base $h(a)$, and their successor equations apply the
same $h$ to the previous equality.  Here and below $I(y,a;h)$ is
an abbreviation for one fixed finite scalar term, not quantification
over functions.

By \eqref{eq:hybrid-factor-map}, another equation induction gives
\begin{equation}
 I(y,3a;H)=3I(y,a;H).
 \label{eq:hybrid-iterator-scaling}
\end{equation}
Its base is $3a=3a$.  In the step the fixed-parameter induction
equation rewrites the new left side as
\[
 H(3I(y,a;H))=3H(I(y,a;H)).
\]
The initializer parameter stays $a$ throughout this induction.
Equation induction also gives $I(y,1;H)=1$, since $H(1)=1$.

Define
\[
 E(x)=I(x,1;u.2u),\qquad
 L(x,y)=N\bigl(I(y,E(x);H)\bigr).
\]
Then the following entire recursion specification is provable:
\begin{equation}
 L(0,y)=1,\qquad L(Sx,0)=0,\qquad L(Sx,Sy)=L(x,y).
 \label{eq:hybrid-comparator}
\end{equation}
The first equation uses $E(0)=1$ and the fixed-point identity for
$H$.  The second uses $E(Sx)=2E(x)$ and $N(2a)=0$.
For the third,
\begin{align*}
 L(Sx,Sy)
 &=N\bigl(I(Sy,2E(x);H)\bigr)\\
 &=N\bigl(I(y,H(2E(x));H)\bigr)\\
 &=N\bigl(I(y,3E(x);H)\bigr)\\
 &=N\bigl(3I(y,E(x);H)\bigr)\\
 &=N\bigl(I(y,E(x);H)\bigr)=L(x,y).
\end{align*}
The successive steps use rotation, $H(2a)=3a$, iterator scaling,
and $N(3a)=N(a)$.  Thus $L$ has the comparison equations at all
arguments internally.  Its standard-model meaning as the indicator
of $x\le y$ is motivation, not a premise of the following argument.

\subsection{Saturated iteration gives predecessor}

We now need only the three equations
\eqref{eq:hybrid-comparator}.  Define
\[
 B_m(u)=C(L(Su,m),u,Su),\qquad
 W(t,m)=I(t,0;u.B_m(u)).
\]
The subscript in $B_m$ records a fixed parameter; $B_m$ is an open
term abbreviation, not a function variable.

Numerical cases on $e$ prove the selector equation
\[
 C(e,Sa,Sb)=S C(e,a,b).
\]
Equivalently one may use equation induction on $e$, whose induction
hypothesis is unused.  Therefore
\begin{align}
 B_{Sm}(Su)
 &=C(L(SSu,Sm),Su,SSu)\notag\\
 &=C(L(Su,m),Su,SSu)=S B_m(u),
 \label{eq:hybrid-conjugacy}\\
 B_{Sm}(0)&=C(L(S0,Sm),0,S0)=C(L(0,m),0,S0)=S0.
 \label{eq:hybrid-conjugacy-base}
\end{align}

Literal-equation induction on $t$, with $m$ fixed, proves
\begin{equation}
 W(St,Sm)=S W(t,m).
 \label{eq:hybrid-saturated-shift}
\end{equation}
The base is \eqref{eq:hybrid-conjugacy-base}.  The successor step is
\begin{align*}
 W(SSt,Sm)
 &=B_{Sm}(W(St,Sm))\\
 &=B_{Sm}(S W(t,m))\\
 &=S B_m(W(t,m))=S W(St,m).
\end{align*}
Only the second line uses the current induction hypothesis, at the
same fixed $m$.

Two subsequent equation inductions give
\begin{equation}
 L(Sn,n)=0,\qquad W(n,n)=n.
 \label{eq:hybrid-diagonals}
\end{equation}
The first uses the comparator recurrence.  The second uses
\eqref{eq:hybrid-saturated-shift} in its successor step.  The
substitution $t=m=n$ is made into that already proved universally
parameterized equation; no parameter of an active induction hypothesis
is changed.  Consequently
\begin{equation}
 W(Sn,n)=B_n(n)=C(L(Sn,n),n,Sn)=n.
 \label{eq:hybrid-extra-step}
\end{equation}

Finally define the finite scalar term
\begin{equation}
 P(m)=I\bigl(m,0;u.C(L(SSu,m),u,Su)\bigr).
 \label{eq:hybrid-predecessor-term}
\end{equation}
Its zero equation is immediate.  To justify its successor equation
without assuming a separate extensionality rule for iterator bodies,
put
\[
 V(t,m)=I\bigl(t,0;u.C(L(SSu,Sm),u,Su)\bigr).
\]
The comparator recurrence identifies this update pointwise with
$B_m(u)$.  Equation induction on $t$, with $m$ fixed, gives
$V(t,m)=W(t,m)$: their base values are $0$, and the successor step
uses the pointwise update equation and then the induction equation.
Now \eqref{eq:hybrid-extra-step} yields
\[
 \boxed{\qquad P(0)=0,\qquad P(Sm)=V(Sm,m)=W(Sm,m)=m.\qquad}
\]

\subsection{Dependency boundary}

All terms in this section are built from $0,S,I$ by finite substitution.
All inductions have a single equation as their invariant.  Numerical
cases and equality transport are the already reconstructed logical
resources.  Successor injectivity may now be obtained by applying $P$
to an assumed equation $Sx=Sy$.  Equality tests, reflection of
quantifier-free formulas, pairing, and the recovery of general
primitive recursion can then be developed in the established order.
They were not used in order to manufacture predecessor.

This argument is a written internal derivation in the stated theory.
The finite implementation checks do not constitute a proof-assistant
formalization of this entire metatheorem.
% Standalone section.  Requires only amsmath and amssymb in the parent document.
% No theorem environments, macros, or bibliography keys are assumed.
% All arithmetic claims in this section are proved here.

\section{From literal-equation induction to full intuitionistic PRA}
\label{app:eq}
\label{sec:eic-completion}\label{app:arithmetic}

This section supplies an unconditional completion theorem. It is adapted
with attribution from the current forward-left draft~\cite{forward}. In the
scalar application used here, every function occurring in the proof is
already a scalar term: predecessor is supplied by
Appendix~\ref{sec:hybrid-bootstrap}, the selector is
$C(n,a,b)=I(n,a;u.b)$, addition is $x+y=I(y,x;u.Su)$, and
truncated subtraction is $d(x,y)=I(y,x;u.Pu)$. All other functions below
are finite compositions of these. Thus this proof does not presuppose
that arbitrary primitive recursion has already been compiled. The QFI
conclusion first applies to the scalar language; the general PR compiler
comes afterward.  Literal-equation
induction, together with the primitive-recursive defining equations,
predecessor coverage, and separation of $0$ and $1$, proves induction for every
quantifier-free formula.  In particular, the proof does not assume a faithful
numerical equality test, stability or decidability of equality, successor
uniqueness, or any commutation law for an iterator.  We construct the equality
test and prove its reflection property for arbitrary elements of the
numerical domain.

\subsection{The exact hypotheses}

Work in ordinary intuitionistic first-order logic with equality and genuine
falsity $\bot$.  The language has one numerical sort and a symbol for each
fixed, finite primitive-recursive description.  Distinct descriptions are
not identified merely because they compute the same standard function.
Let $\mathsf{PRE}$ contain the defining equations for zero, successor,
projections, composition, and primitive recursion.  In particular, choose
descriptions $P$ and $C$ with equations
\begin{align}
 P(0)&=0, & P(Sx)&=x, \label{eq:eic-predecessor}\\
 C(0,a,b)&=a, & C(Sx,a,b)&=b. \label{eq:eic-selector}
\end{align}
Write $1=S0$.  In addition to $\mathsf{PRE}$, assume
\begin{align}
 \mathsf{Cov}:&\qquad
       \forall x\,\bigl(x=0\lor x=S(Px)\bigr),
       \label{eq:eic-coverage}\\
 \mathsf{Sep}:&\qquad 0=1\longrightarrow\bot,
       \label{eq:eic-separation}
\end{align}
and induction for each \emph{literal equation} $E(n,\vec p)$:
\begin{equation}
 \forall\vec p\,
 \left[
   \left(E(0,\vec p)\land
       \forall n\bigl(E(n,\vec p)\longrightarrow E(Sn,\vec p)\bigr)
   \right)
   \longrightarrow\forall n\,E(n,\vec p)
 \right].
 \label{eq:eic-induction}
\end{equation}
Here the schematic expression $E$ is a single equality of numerical
terms; the displayed universal quantifiers range only over numbers.
Every description and every instance of the schema is finite external
syntax.

The axiom formulation \eqref{eq:eic-induction} and the usual contextual
induction rule are interchangeable.  For example, if
\[
 \Gamma\vdash E(0,\vec p),
 \qquad
 \Gamma,E(z,\vec p)\vdash E(Sz,\vec p),
\]
where $z$ is absent from $\Gamma$ and the passive parameters, discharge the
second premise's equation and universally introduce $z$.  The axiom then
gives $\Gamma\vdash\forall n E(n,\vec p)$, which can be instantiated at any
target term.  Conversely, assume the antecedent of the axiom and apply the
contextual rule with fresh variables.  This uses the ordinary intuitionistic
logical rules and permits arbitrary, including quantified, formulas in
$\Gamma$.

\subsection{Constructor separation and two elementary tests}

Define the sign description and the countdown description by
\begin{equation}
 \operatorname{sg}(u)=C(u,0,1),
 \qquad
 K(u)=C(u,1,Pu).
 \label{eq:eic-sign-countdown}
\end{equation}
Their defining equations imply
\begin{align}
 \operatorname{sg}(0)&=0,
 &\operatorname{sg}(Sx)&=1,\label{eq:eic-sign-computation}\\
 K(0)&=1,
 &K(Sx)&=x.\label{eq:eic-countdown-computation}
\end{align}

\paragraph{Constructor laws.}
Successor is injective: applying $P$ to $Sx=Sy$ gives $x=y$ by
\eqref{eq:eic-predecessor}.  Successor never equals zero.  Indeed,
$Sx=0$ gives $1=0$ by applying $\operatorname{sg}$ and using
\eqref{eq:eic-sign-computation}; symmetry and
\eqref{eq:eic-separation} give $\bot$.  Consequently $1\ne0$, where
$A\ne B$ abbreviates $A=B\to\bot$.  None of these arguments uses
induction.

\paragraph{The sign test.}
For every $u$,
\begin{align}
 &\operatorname{sg}(u)=0\lor\operatorname{sg}(u)=1,
       \label{eq:eic-sign-range}\\
 &\operatorname{sg}(u)=0\quad\longleftrightarrow\quad u=0,
       \label{eq:eic-sign-zero}\\
 &\operatorname{sg}(u)=1\quad\longleftrightarrow\quad u\ne0.
       \label{eq:eic-sign-one}
\end{align}
To prove the range assertion, split \eqref{eq:eic-coverage}.  The branch
$u=0$ gives sign $0$, and the branch $u=S(Pu)$ gives sign $1$.
The implication $u=0\to\operatorname{sg}(u)=0$ is computation.
For its converse, the successor branch of coverage would make
$\operatorname{sg}(u)=0$ imply $1=0$, a contradiction; in that branch
explosion gives $u=0$ as well.  Similarly, if
$\operatorname{sg}(u)=1$ and $u=0$, computation gives $0=1$.
Conversely, $u\ne0$ excludes the zero branch of coverage, and the
successor branch computes the sign as $1$.

Coverage and \eqref{eq:eic-selector} also give, for all $u,v$,
\begin{equation}
 C(u,v,v)=v.
 \label{eq:eic-selector-idempotence}
\end{equation}
Both coverage branches compute to $v$, so this identity again needs no
induction.

\paragraph{The countdown test and the notation $W$.}
For every $u$,
\begin{equation}
 K(u)=0\quad\longleftrightarrow\quad u=1.
 \label{eq:eic-countdown-zero}
\end{equation}
If $u=1=S0$, \eqref{eq:eic-countdown-computation} gives $K(u)=0$.
Conversely, assume $K(u)=0$ and split coverage.  If $u=0$, then
$K(u)=1$, contradicting separation, and explosion gives $u=1$.
If $u=S(Pu)$, computation gives $K(u)=Pu$; hence $Pu=0$ and
$u=S0=1$.

If a presentation uses a description named $W$ with defining equations
\[
 W(0)=1,\qquad W(Sx)=x,
\]
then it proves $K(u)=W(u)$ for every $u$ already from coverage.
In the zero branch both sides are $1$; in the branch $u=S(Pu)$ both
sides are $Pu$.  Thus $K$ and $W$ are provably equal on the whole domain,
even if they are distinct descriptions.  This equality does not appeal to
standard-model agreement or to uniqueness of an arbitrary recurrence.

\subsection{Four elementary uses of equation induction}

Choose descriptions for addition and truncated subtraction, with recursion
on the second argument:
\begin{align}
 x+0&=x, &x+Sy&=S(x+y),\label{eq:eic-addition}\\
 d(x,0)&=x, &d(x,Sy)&=P(d(x,y)).\label{eq:eic-subtraction}
\end{align}
Using these equations and \eqref{eq:eic-induction}, we now derive
\begin{align}
 0+y&=y,\label{eq:eic-left-zero}\\
 (Sx)+y&=S(x+y),\label{eq:eic-left-successor}\\
 d(Sx,Sy)&=d(x,y),\label{eq:eic-cancel-successors}\\
 d(Sa,a)&=1.\label{eq:eic-diagonal}
\end{align}
Each induction below has a single equation as its induction formula.

For \eqref{eq:eic-left-zero}, induct on $y$.  The base is $0+0=0$.
The step is $0+Sy=S(0+y)=Sy$.

For \eqref{eq:eic-left-successor}, hold $x$ fixed and induct on $y$.
At zero both sides are $Sx$.  At a successor the induction hypothesis
gives
\[
 (Sx)+Sy=S\bigl((Sx)+y\bigr)
       =S\bigl(S(x+y)\bigr)=S(x+Sy).
\]
Combining \eqref{eq:eic-left-successor} with the defining equation for
addition gives the useful rearrangement
\begin{equation}
 (Sy)+r=y+Sr.
 \label{eq:eic-shift-successor}
\end{equation}
Indeed, both sides are $S(y+r)$.

For \eqref{eq:eic-cancel-successors}, hold $x$ fixed and induct on $y$.
The base computes as
\[
 d(Sx,S0)=P(d(Sx,0))=P(Sx)=x=d(x,0).
\]
Applying $P$ to the induction equation gives the step:
\[
 d(Sx,S(Sy))=P(d(Sx,Sy))=P(d(x,y))=d(x,Sy).
\]

Finally, induct on $a$ for \eqref{eq:eic-diagonal}.  The base is
$d(S0,0)=S0=1$.  At a successor,
\[
 d(S(Sa),Sa)=d(Sa,a)=1,
\]
using \eqref{eq:eic-cancel-successors} and the induction hypothesis.

\subsection{An equation invariant that reflects actual equality}

The remaining arithmetic ingredient is the following equation, valid for
all $x,y$:
\begin{equation}
 C\bigl(\operatorname{sg}(d(x,y)),\ x,\ y+d(x,y)\bigr)=x.
 \label{eq:eic-residual-invariant}
\end{equation}
Its meaning is that a nonzero subtraction residual reconstructs the
original argument by addition.  The selector expresses this assertion as
one literal equation, so equation induction applies directly.

Hold $x$ fixed and induct on $y$.  At $y=0$, subtraction gives residual
$x$, and \eqref{eq:eic-left-zero} gives $0+x=x$.  The desired equation
is consequently $C(\operatorname{sg}(x),x,x)=x$, which is
\eqref{eq:eic-selector-idempotence}.

For the step, assume \eqref{eq:eic-residual-invariant} at $y$ and
abbreviate
\[
 r=d(x,y),\qquad r'=Pr=d(x,Sy).
\]
These are local term abbreviations, not additional function symbols.
Split coverage for $r'$.

If $r'=0$, then $\operatorname{sg}(r')=0$, and the successor instance
of \eqref{eq:eic-residual-invariant} reduces immediately to $x=x$.

In the other branch, $r'=S(Pr')$, so constructor separation gives
$r'\ne0$.  Split coverage for $r$.  Its zero branch would imply
$r'=P0=0$, a contradiction, from which the desired step equation follows.
In its successor branch,
\[
 r=S(Pr)=Sr',\qquad \operatorname{sg}(r)=1.
\]
The induction equation therefore computes to $y+r=x$.  Meanwhile,
\eqref{eq:eic-shift-successor} gives
\[
 (Sy)+r'=y+Sr'=y+r=x.
\]
Since $r'$ is a successor, $\operatorname{sg}(r')=1$, so this is
exactly the branch selected in the successor instance of
\eqref{eq:eic-residual-invariant}.  This completes the equation induction.
All case distinctions came from coverage; no induction on an implication
or a disjunction was used.

\paragraph{Reflection at residual $1$.}
We have, for arbitrary $a,b$,
\begin{equation}
 d(Sa,b)=1\quad\longleftrightarrow\quad a=b.
 \label{eq:eic-residual-reflection}
\end{equation}
For the forward direction, substitute $x=Sa$, $y=b$ into
\eqref{eq:eic-residual-invariant}.  If $d(Sa,b)=1$, the sign is $1$
and the selector returns its second branch, giving
\[
 b+1=Sa.
\]
The defining equation for addition gives $b+1=S(b+0)=Sb$.  Thus
$Sb=Sa$, and applying $P$ gives $b=a$, hence $a=b$.
Conversely, $a=b$ and \eqref{eq:eic-diagonal} give
$d(Sa,b)=d(Sa,a)=1$.

Define the equality characteristic by the fixed description
\begin{equation}
 e(a,b)=\operatorname{sg}\bigl(K(d(Sa,b))\bigr).
 \label{eq:eic-equality-characteristic}
\end{equation}
The sign, countdown, and residual reflection lemmas give
\begin{align}
 e(a,b)=0&\quad\longleftrightarrow\quad a=b,
       \label{eq:eic-equality-reflection}\\
 e(a,b)=0&\ \lor\ e(a,b)=1.
       \label{eq:eic-equality-range}
\end{align}
In detail, $e(a,b)=0$ is equivalent to $K(d(Sa,b))=0$ by
\eqref{eq:eic-sign-zero}, which is equivalent to $d(Sa,b)=1$ by
\eqref{eq:eic-countdown-zero}, and then to $a=b$ by
\eqref{eq:eic-residual-reflection}.  The range assertion is
\eqref{eq:eic-sign-range} applied to $K(d(Sa,b))$.

The reflection property concerns the logic's actual equality throughout
the domain.  It was not assumed as part of a numerical comparison symbol,
and its proof did not restrict $a$ or $b$ to standard numerals.
It also yields decidable equality.  In the $e(a,b)=0$ branch,
\eqref{eq:eic-equality-reflection} gives $a=b$.  In the $e(a,b)=1$
branch, an assumption $a=b$ would give $e(a,b)=0$ and therefore $0=1$,
contradicting separation.  Hence
\begin{equation}
 \forall a\,\forall b\,\bigl(a=b\lor(a=b\to\bot)\bigr)
 \label{eq:eic-equality-decidable}
\end{equation}
is a theorem, not an extra hypothesis.

\subsection{Compilation of all quantifier-free formulas}

For each fixed quantifier-free formula $A(\vec p)$, construct a term
$b_A(\vec p)$ by external recursion on its finite syntax.  Take equality,
$\bot$, $\land$, $\lor$, and $\to$ as the primitive formula constructors;
negation is implication into $\bot$.  Set
\begin{align}
 b_{s=t}&=e(s,t), & b_{\bot}&=1,\label{eq:eic-boolean-atoms}\\
 b_{A\land D}&=C(b_A,b_D,1),\label{eq:eic-boolean-and}\\
 b_{A\lor D}&=C(b_A,0,b_D),\label{eq:eic-boolean-or}\\
 b_{A\to D}&=C(b_A,b_D,0).\label{eq:eic-boolean-implies}
\end{align}
Every term here is a finite composition of fixed primitive-recursive
descriptions; if desired, projections and composition give it a single
description symbol.  No internal quantification over formulas or programs
is involved.

An external induction on the syntax of $A$ proves, within the arithmetic
theory,
\begin{equation}
 \bigl(b_A=0\lor b_A=1\bigr)
 \quad\text{and}\quad
 \bigl(b_A=0\longleftrightarrow A\bigr).
 \label{eq:eic-boolean-correctness}
\end{equation}
Here are the full connective checks.  In every case $b_A=1$ implies
$\neg A$: if $A$ held, the already established implication
$A\to b_A=0$ would contradict separation.

\paragraph{Atomic formulas and falsity.}
For an equation, both claims are
\eqref{eq:eic-equality-reflection}--\eqref{eq:eic-equality-range}.
For falsity the range assertion is $1=0\lor1=1$, with its right disjunct
proved by reflexivity.  Also $1=0\to\bot$ follows from separation and
symmetry, while $\bot\to1=0$ follows by explosion.  Thus $1=0$ is
equivalent to $\bot$, as required.

\paragraph{Conjunction.}
Split $b_A=0\lor b_A=1$.  In the zero branch, $A$ holds and
$b_{A\land D}=b_D$.  Its range follows from the range of $b_D$,
and its being zero is equivalent to $D$, hence to $A\land D$ under
the known $A$.  In the one branch, $\neg A$ holds and
$b_{A\land D}=1$.  Its range is immediate.  If it were zero,
separation would give $\bot$, hence $A\land D$; conversely an
assumption $A\land D$ contradicts $\neg A$, and explosion gives
$b_{A\land D}=0$.

\paragraph{Disjunction.}
In the branch $b_A=0$, the formula $A$ holds and
$b_{A\lor D}=0$.  The range is immediate, and both $b_{A\lor D}=0$
and $A\lor D$ hold, proving their equivalence.  In the branch
$b_A=1$, we have $\neg A$ and $b_{A\lor D}=b_D$.  Its range follows
from the range of $b_D$.  If its value is zero, then $D$ and hence
$A\lor D$.  Conversely, from $A\lor D$, the $A$ alternative
contradicts $\neg A$ and yields $b_D=0$ by explosion, while the
$D$ alternative yields $b_D=0$ directly.  This proves the required
equivalence in both alternatives.

\paragraph{Implication.}
In the branch $b_A=0$, the formula $A$ holds and
$b_{A\to D}=b_D$.  Its range follows from the range of $b_D$.
If its value is zero, then $D$, which gives $A\to D$.  Conversely,
$A\to D$ and the known $A$ give $D$, hence $b_D=0$.
In the branch $b_A=1$, we have $\neg A$ and
$b_{A\to D}=0$.  Its range is immediate.  To prove $A\to D$,
assume $A$; this contradicts $\neg A$, and explosion gives $D$.
Thus both $b_{A\to D}=0$ and $A\to D$ hold in this branch, and
their equivalence follows.  This argument uses the proved two-value
range of the finite characteristic.  It does not use excluded middle
for arbitrary formulas or change the surrounding intuitionistic logic.

This completes the external syntactic induction proving
\eqref{eq:eic-boolean-correctness} for every quantifier-free formula.

\subsection{Recovery of induction and the full theory}

Let $A(n,\vec p)$ be any fixed quantifier-free formula.  Suppose an
arbitrary retained context $\Gamma$ proves its base and its induction
step:
\[
 \Gamma\vdash A(0,\vec p),
 \qquad
 \Gamma,A(z,\vec p)\vdash A(Sz,\vec p),
\]
with the usual freshness condition on $z$.  By
\eqref{eq:eic-boolean-correctness}, the base gives
$\Gamma\vdash b_A(0,\vec p)=0$.  Under the additional assumption
$b_A(z,\vec p)=0$, the same equivalence gives $A(z,\vec p)$,
the given step gives $A(Sz,\vec p)$, and the equivalence gives
$b_A(Sz,\vec p)=0$.  Apply literal-equation induction to
$b_A(n,\vec p)=0$ and then apply its reflection implication.
The result is $\Gamma\vdash A(t,\vec p)$ for every target term $t$.
Thus the theory proves the full contextual quantifier-free induction
rule, including in arbitrary quantified contexts.

\paragraph{Completion theorem.}
In the common language of fixed primitive-recursive descriptions, the
theory
\[
 \mathsf{PRE}+\mathsf{Cov}+\mathsf{Sep}+\mathsf{EI}
\]
has exactly the same sequents as ordinary first-order intuitionistic
primitive recursive arithmetic with all quantifier-free induction
instances.  The forward simulation follows from the recovered
quantifier-free induction, successor injectivity, and successor separation;
all primitive-recursive defining equations and all logical rules are
shared.  For the reverse simulation, literal-equation induction is an
instance of quantifier-free induction, and successor separation gives
$\mathsf{Sep}$.  Quantifier-free induction proves $\mathsf{Cov}$:
at zero choose its left disjunct, and at a successor use
$Sx=S(P(Sx))$, which follows from $P(Sx)=x$, for its right disjunct.

Both interpretations therefore use the full domain, defined by $x=x$,
and preserve every function symbol, equality, logical connective, and
quantifier.  Both composites have the identity isomorphism, with graph
$x=y$.  The theories are consequently identity-domain bi-interpretable
and equiconsistent.  This completion theorem is unconditional once the
literal-equation induction rule has been obtained; it has no dependency
on a classical successor-uniqueness reduction.
\section{Pairing, primitive recursion, and internal numerical composites}
\label{app:pr}
\label{sec:iterator-coding}
This appendix is adapted with attribution from the progress-closure
arithmetic development~\cite{progress}. Its prerequisites have now been
proved: scalar predecessor, arithmetic reflection, and full quantifier-free
induction. Pairing is a derived proof tool; it is not part of any checker.

\subsection{Pairing from a bounded square-root search}

Once QFI is available, a familiar square-based injection gives a
particularly direct compiler. Write $a\le b$ for $d(a,b)=0$.
The usual linear-order and semiring identities follow by fixed
quantifier-free inductions from the successor, addition, and monus
equations. Multiplication itself is the scalar term
\[
                        a\cdot b=I(b,0;u.u+a).
\]
Only primitive-recursive arithmetic is involved in the following
bounded root search; there is no minimization over an unbounded domain.

For a fixed parameter $n$, put
\[
 U_n(v)=C(d((Sv)^2,n),Sv,v),\qquad
 R_n(i)=I(i,0;v.U_n(v)),\qquad q(n)=R_n(n).
\]
The update increments $v$ exactly when $(v+1)^2\le n$, and
otherwise retains it. QFI in $i$, with $n$ fixed, proves
\[
 R_n(i)\le i,\qquad R_n(i)^2\le n,\qquad
 R_n(i)=i\ \vee\ n<(S(R_n(i)))^2.              \tag{Root invariant}
\]
The base has $R_n(0)=0$. If the update increments, its test
contradicts the strict-upper-bound alternative in the third invariant.
Thus the old state equals $i$, and the new state equals $Si$;
its square bound is the increment test. If the update does not increment, the
strict upper square bound holds and the state remains fixed.
An already established strict upper square bound forces every
subsequent update to remain fixed. These are exactly the cases
needed for the three displayed invariants.

At $i=n$, either the strict upper square bound already holds or
$R_n(n)=n$, in which case $n<(n+1)^2$ supplies it. Thus
\[
                        q(n)^2\le n<(q(n)+1)^2.           \tag{Root}
\]
Monotonicity of squares shows uniqueness of a number satisfying
Root. All formulas just used are quantifier free in the scalar
language, with fixed $n$ during the induction.

Now define
\[
 \pi(a,b)=(a+b)^2+a,\qquad
 L(n)=d(n,q(n)^2),\qquad H(n)=d(q(n),L(n)).
\]
For $k=a+b$, one has $k^2\le k^2+a<(k+1)^2$ since $a\le k$.
Root uniqueness gives $q(\pi(a,b))=a+b$, and therefore
\[
                    L(\pi(a,b))=a,\qquad H(\pi(a,b))=b.
\]
Surjectivity of $\pi$ is unnecessary. The original numerical domain
is never replaced by the range of this coding; pairing is used only
to store the finite state of a fixed primitive recursion. This
standard square pairing is a convenient reconstruction here, not
a claimed new arithmetic basis.

\subsection{Compiling each primitive recursion}

Suppose the fixed target description is
\[
 f(0,\vec p)=g(\vec p),\qquad
 f(Sn,\vec p)=h(n,f(n,\vec p),\vec p).
\]
Assume by external structural induction that the earlier $g,h$
already have scalar translations $g^*,h^*$. Define
\[
 W(u,\vec p)=\pi(SL(u),h^*(L(u),H(u),\vec p)),
\]
\[
 V(n,\vec p)=I(n,\pi(0,g^*(\vec p));u.W(u,\vec p)),
 \qquad f^*(n,\vec p)=H(V(n,\vec p)).
\]
The update has only its state variable bound, with every $p_i$
fixed. Equation induction proves $L(V(n,\vec p))=n$.
The two projection equations and scalar computation then prove
both defining equations for $f^*$ internally. Initial functions
and composition translate directly. This completes the translation
of every fixed finite PR description without adding a universal
evaluator.

\subsection{Why both numerical round trips are internal}

Two fixed scalar iterators with the same seed and provably equal
updates are internally equal at all counters: use literal-equation
induction comparing their values. Their base agrees, and their
successor equations plus congruence prove the step. The same
uniqueness argument works for ordinary primitive recursion.

For an ordinary PR description, induct externally on its finite
construction. The double translation satisfies its original
defining equations, so internal equation induction identifies it
with the original operation. For a native scalar description,
earlier translated updates agree by the external induction
hypothesis; scalar uniqueness identifies the double-translated
iterator. These proofs hold at arbitrary internal counters.

Thus the domain relation $x=y$ is provably a total bijection
preserving every fixed numerical operation in both composites.
The formula constructions and round trips in
Section~\ref{sec:choice-introduction} give the corresponding formula
composites. Both interpretations preserve falsity, taking native
$0=1$ to an ordinary contradiction and ordinary $\bot$ to native
$0=1$. Finite proof translations therefore establish equiconsistency.

\clearpage
\section{Complete raw root-free checker}
The following listing is the actual standalone selected implementation,
including parsing, validation, inference, and its public interface.
\begin{lstlisting}[language=Python]
"""Raw clausal iPRA with computational equality.

Terms: 0, negative free levels, positive bound indices, [0,t] (S),
[n,a,h] (I). Formulas: [s,t] (equality), [B] (D).
Contexts contain (antecedent, one-bound-variable body) clauses.
Input: {'goal': a closed formula, 'proof': a finite proof tree}.
"""
import json
import sys


def fail():
    raise ValueError('invalid certificate')


def term(a, scope, depth=0):
    if type(a) is int:
        return -scope <= a <= depth
    return (type(a) is list and len(a) in (2, 3)
            and (len(a) == 3 or a[0] == 0)
            and all(term(t, scope, depth + (len(a) == 3 and i == 2))
                    for i, t in enumerate(a)))


def formula(a, scope, depth=0):
    return (type(a) is list
            and ((len(a) == 2 and all(term(t, scope, depth) for t in a))
                 or (len(a) == 1 and formula(a[0], scope, depth + 2))))


def opening(a, value, depth=0):
    if type(a) is int:
        return value if a == depth + 1 else a - (a > depth + 1)
    size = len(a)
    return [opening(t, value, depth + (2 if size == 1 else int(size == 3 and i == 2)))
            for i, t in enumerate(a)]


def derive(proof, context, scope):
    if type(proof) is not list or not proof or type(proof[0]) is not str:
        fail()
    tag = proof[0]
    if tag == 'u' and len(proof) == 4:
        index, witness, premise = proof[1:]
        if (type(index) is not int or not 0 <= index < len(context)
                or not term(witness, scope)):
            fail()
        antecedent, body = context[index]
        if derive(premise, context, scope) != antecedent:
            fail()
        return opening(body, witness)
    if tag == 'r' and len(proof) == 2:
        a = proof[1]
        if not term(a, scope):
            fail()
        return [a, a]
    if tag == 'v' and len(proof) == 3:
        s, t = derive(proof[1], context, scope)
        b, r = derive(proof[2], context, scope)
        c, a, h = b
        if c != 0:
            if len(c) != 2:
                fail()
            a, t = [c[1], a, h], opening(h, t)
        if a != s:
            fail()
        return [r, t]
    if tag in ('e', 'i') and len(proof) == 4:
        goal, data, branch = proof[1:]
        if not formula(goal, scope):
            fail()
        fresh = -scope - 1
        if tag == 'e':
            paired = derive(data, context, scope)
            if len(paired) != 1:
                fail()
            body = opening(paired[0], fresh, 1)
            clause = (opening(body, 0), body)
            target = goal
        else:
            if len(goal) != 1 or not term(data, scope):
                fail()
            body = opening(goal[0], data, 1)
            eta = fresh if len(goal[0]) == 2 and data == 0 else 0
            clause = ([0, 0], opening(body, eta))
            target = opening(body, [0, fresh])
        if derive(branch, context + [clause], scope + 1) != target:
            fail()
        return goal
    fail()


def verify(data):
    if type(data) is not dict or set(data) != {'goal', 'proof'}:
        fail()
    if not formula(data['goal'], 0):
        fail()
    if derive(data['proof'], [], 0) != data['goal']:
        fail()
    return True


if __name__ == '__main__':
    try:
        verify(json.load(sys.stdin))
    except Exception:
        print('REJECT')
        sys.exit(1)
    print('ACCEPT')
\end{lstlisting}

\clearpage
\section{Complete zero-normal companion checker}
This separate implementation includes canonical syntax validation,
normalizing opening, and the fixed-zero-base introduction. It uses the
same external JSON wrapper.
\begin{lstlisting}[language=Python]
"""Zero-normal clausal iPRA with uniform evaluation and zero-base introduction.

Terms: 0, negative free levels, positive bound indices, [0,t] (S),
[n,a,h] (I), n not literally 0. Formulas: [s,t] (equality), [B] (D).
Contexts contain (antecedent, one-bound-variable body) clauses.
Input: {'goal': a closed formula, 'proof': a finite proof tree}.
"""
import json
import sys


def fail():
    raise ValueError('invalid certificate')


def term(a, scope, depth=0):
    if type(a) is int:
        return -scope <= a <= depth
    return (type(a) is list and len(a) in (2, 3)
            and (len(a) == 2) == (a[0] == 0)
            and all(term(t, scope, depth + (len(a) == 3 and i == 2))
                    for i, t in enumerate(a)))


def formula(a, scope, depth=0):
    return (type(a) is list
            and ((len(a) == 2 and all(term(t, scope, depth) for t in a))
                 or (len(a) == 1 and formula(a[0], scope, depth + 2))))


def opening(a, value, depth=0):
    if type(a) is int:
        return value if a == depth + 1 else a - (a > depth + 1)
    size = len(a)
    a = [opening(t, value, depth + (2 if size == 1 else int(size == 3 and i == 2)))
         for i, t in enumerate(a)]
    return a[1] if size == 3 and a[0] == 0 else a


def derive(proof, context, scope):
    if type(proof) is not list or not proof or type(proof[0]) is not str:
        fail()
    tag = proof[0]
    if tag == 'u' and len(proof) == 4:
        index, witness, premise = proof[1:]
        if (type(index) is not int or not 0 <= index < len(context)
                or not term(witness, scope)):
            fail()
        antecedent, body = context[index]
        if derive(premise, context, scope) != antecedent:
            fail()
        return opening(body, witness)
    if tag == 'r' and len(proof) == 2:
        a = proof[1]
        if not term(a, scope):
            fail()
        return [a, a]
    if tag == 'v' and len(proof) == 3:
        s, t = derive(proof[1], context, scope)
        b, r = derive(proof[2], context, scope)
        c, a, h = b
        if len(c) != 2 or s != (a if c[1] == 0 else [c[1], a, h]):
            fail()
        return [r, opening(h, t)]
    if tag in ('e', 'i') and len(proof) == 4:
        goal, data, branch = proof[1:]
        if not formula(goal, scope):
            fail()
        fresh = -scope - 1
        if tag == 'e':
            paired = derive(data, context, scope)
            if len(paired) != 1:
                fail()
            body = opening(paired[0], fresh, 1)
            clause = (opening(body, 0), body)
            target = goal
        else:
            if len(goal) != 1 or not term(data, scope):
                fail()
            body = opening(goal[0], data, 1)
            base = opening(body, 0)
            clause = ([0, 0], opening(body, fresh) if base == [0, 0] else base)
            target = opening(body, [0, fresh])
        if derive(branch, context + [clause], scope + 1) != target:
            fail()
        return goal
    fail()


def verify(data):
    if type(data) is not dict or set(data) != {'goal', 'proof'}:
        fail()
    if not formula(data['goal'], 0):
        fail()
    if derive(data['proof'], [], 0) != data['goal']:
        fail()
    return True


if __name__ == '__main__':
    try:
        verify(json.load(sys.stdin))
    except Exception:
        print('REJECT')
        sys.exit(1)
    print('ACCEPT')
\end{lstlisting}

\end{document}
